% =====================================================================
%  Chapitre 6 : systèmes multi-qubits et intrication
%  Fichier importé par main.tex avec % =====================================================================
%  Chapitre 6 : systèmes multi-qubits et intrication
%  Fichier importé par main.tex avec % =====================================================================
%  Chapitre 6 : systèmes multi-qubits et intrication
%  Fichier importé par main.tex avec % =====================================================================
%  Chapitre 6 : systèmes multi-qubits et intrication
%  Fichier importé par main.tex avec \include{chapitre6}.
%  Comme chapitre4.tex et chapitre5.tex, il ouvre une NOUVELLE SECTION
%  (numéro 4). Pour le rattacher à la section précédente, remplacer
%  \section par \subsection et chaque \subsection par \subsubsection.
%  Pas de préambule : environnements (definition, resultat, remarque),
%  macros (\ii, \ee, \dd, \ket, \bra, \braket) et couleurs (insarouge,
%  insagris) sont définis dans main.tex.
% =====================================================================

% Macros de Dirac (déjà définies dans main.tex récent ; ce bloc évite l'erreur
% « Undefined control sequence \ket » si main.tex est une ancienne version).
\providecommand{\ket}[1]{|#1\rangle}
\providecommand{\bra}[1]{\langle #1|}
\providecommand{\braket}[2]{\langle #1|#2\rangle}

% =====================================================================
\section{Systèmes multi-qubits et intrication}\label{sec:multi}
% =====================================================================

Un ordinateur quantique utilise plusieurs qubits. Cette section décrit l'espace des états
de deux qubits (le produit tensoriel), la mesure, la distinction entre états séparables
et états intriqués, puis les opérateurs qui agissent sur deux qubits. Tous les calculs
sont détaillés ; les qubits sont notés $A$ et $B$.

% ---------------------------------------------------------------------
\subsection{L'espace de deux qubits}\label{sec:tenseur}
% ---------------------------------------------------------------------

Chaque qubit vit dans $\mathcal{H}^2=\mathbb{C}^2$. L'espace du système $(A,B)$ est le
\emph{produit tensoriel}
\[
  \mathcal{H}^2\otimes\mathcal{H}^2=\mathcal{H}^4=\mathbb{C}^4 ,
\]
de dimension $2\times2=4$. Pour $n$ qubits la dimension est $2^n$ :

\begin{center}
\renewcommand{\arraystretch}{1.3}
\begin{tabular}{@{}lccccc@{}}
\toprule
Nombre de qubits $n$ & $1$ & $2$ & $3$ & $10$ & $50$ \\
\midrule
Dimension $2^n$ & $2$ & $4$ & $8$ & $1024$ & $\approx1{,}1\times10^{15}$ \\
\bottomrule
\end{tabular}
\end{center}

\begin{definition}{Produit tensoriel de deux vecteurs}
\[
  \begin{pmatrix}a_1\\ a_2\end{pmatrix}\otimes\begin{pmatrix}b_1\\ b_2\end{pmatrix}
  =\begin{pmatrix}a_1\begin{pmatrix}b_1\\ b_2\end{pmatrix}\\[2ex]
                  a_2\begin{pmatrix}b_1\\ b_2\end{pmatrix}\end{pmatrix}
  =\begin{pmatrix}a_1b_1\\ a_1b_2\\ a_2b_1\\ a_2b_2\end{pmatrix}.
\]
Chaque composante du premier vecteur multiplie le second vecteur en entier. On note
$\ket{ij}_{AB}=\ket{i}_A\otimes\ket{j}_B$ : le premier chiffre est l'état de $A$, le second
celui de $B$.
\end{definition}

\subsubsection{Les quatre états de base}

On développe chacun des quatre produits avec $\ket{0}=\begin{pmatrix}1\\ 0\end{pmatrix}$
et $\ket{1}=\begin{pmatrix}0\\ 1\end{pmatrix}$ :
\begin{align*}
  \ket{0}\otimes\ket{0}&=\begin{pmatrix}1\\ 0\end{pmatrix}\otimes\begin{pmatrix}1\\ 0\end{pmatrix}
    =\begin{pmatrix}1\cdot\begin{pmatrix}1\\ 0\end{pmatrix}\\[1.5ex]0\cdot\begin{pmatrix}1\\ 0\end{pmatrix}\end{pmatrix}
    =\begin{pmatrix}1\\ 0\\ 0\\ 0\end{pmatrix}=\ket{00},\\[1.5ex]
  \ket{0}\otimes\ket{1}&=\begin{pmatrix}1\\ 0\end{pmatrix}\otimes\begin{pmatrix}0\\ 1\end{pmatrix}
    =\begin{pmatrix}1\cdot\begin{pmatrix}0\\ 1\end{pmatrix}\\[1.5ex]0\cdot\begin{pmatrix}0\\ 1\end{pmatrix}\end{pmatrix}
    =\begin{pmatrix}0\\ 1\\ 0\\ 0\end{pmatrix}=\ket{01},\\[1.5ex]
  \ket{1}\otimes\ket{0}&=\begin{pmatrix}0\\ 1\end{pmatrix}\otimes\begin{pmatrix}1\\ 0\end{pmatrix}
    =\begin{pmatrix}0\cdot\begin{pmatrix}1\\ 0\end{pmatrix}\\[1.5ex]1\cdot\begin{pmatrix}1\\ 0\end{pmatrix}\end{pmatrix}
    =\begin{pmatrix}0\\ 0\\ 1\\ 0\end{pmatrix}=\ket{10},\\[1.5ex]
  \ket{1}\otimes\ket{1}&=\begin{pmatrix}0\\ 1\end{pmatrix}\otimes\begin{pmatrix}0\\ 1\end{pmatrix}
    =\begin{pmatrix}0\cdot\begin{pmatrix}0\\ 1\end{pmatrix}\\[1.5ex]1\cdot\begin{pmatrix}0\\ 1\end{pmatrix}\end{pmatrix}
    =\begin{pmatrix}0\\ 0\\ 0\\ 1\end{pmatrix}=\ket{11}.
\end{align*}
Le vecteur $\ket{ij}$ a son $1$ en position $2i+j$ (en comptant à partir de $0$) : c'est
la lecture de $ij$ en binaire. Les quatre vecteurs forment la base de calcul de
$\mathcal{H}^4$.

\subsubsection{Produit scalaire et orthonormalité}

\begin{resultat}{Produit scalaire d'un produit tensoriel}
\[
  \big(\bra{a}\otimes\bra{b}\big)\big(\ket{c}\otimes\ket{d}\big)=\braket{a}{c}\,\braket{b}{d}.
\]
\end{resultat}

\emph{Démonstration.} Les composantes de $\ket{a}\otimes\ket{b}$ sont $a_ib_j$, celles de
$\ket{c}\otimes\ket{d}$ sont $c_id_j$. Le produit scalaire est
$\sum_{i,j}(a_ib_j)^*c_id_j=\big(\sum_ia_i^*c_i\big)\big(\sum_jb_j^*d_j\big)=\braket{a}{c}\braket{b}{d}$.
\hfill$\square$

On en déduit $\braket{ij}{kl}=\braket{i}{k}\braket{j}{l}$. Les seize produits scalaires
des états de base sont :
\begin{center}
\renewcommand{\arraystretch}{1.3}
\begin{tabular}{@{}c|cccc@{}}
 & $\ket{00}$ & $\ket{01}$ & $\ket{10}$ & $\ket{11}$ \\
\hline
$\bra{00}$ & $1\cdot1=1$ & $1\cdot0=0$ & $0\cdot1=0$ & $0\cdot0=0$ \\
$\bra{01}$ & $1\cdot0=0$ & $1\cdot1=1$ & $0\cdot0=0$ & $0\cdot1=0$ \\
$\bra{10}$ & $0\cdot1=0$ & $0\cdot0=0$ & $1\cdot1=1$ & $1\cdot0=0$ \\
$\bra{11}$ & $0\cdot0=0$ & $0\cdot1=0$ & $1\cdot0=0$ & $1\cdot1=1$ \\
\end{tabular}
\end{center}
(chaque case est $\braket{i}{k}\cdot\braket{j}{l}$). On lit $\braket{ij}{kl}=\delta_{ik}\delta_{jl}$ :
la base est \emph{orthonormée}.

\subsubsection{Produit tensoriel d'états quelconques}

Le produit tensoriel est linéaire en chaque facteur :
$(\lambda\ket{a}+\mu\ket{a'})\otimes\ket{b}=\lambda\ket{a}\otimes\ket{b}+\mu\ket{a'}\otimes\ket{b}$,
et de même pour le second facteur. Pour $\alpha\ket{0}+\beta\ket{1}$ et
$\gamma\ket{0}+\delta\ket{1}$ :
\begin{align*}
  \big(\alpha\ket{0}+\beta\ket{1}\big)\otimes\big(\gamma\ket{0}+\delta\ket{1}\big)
  &=\alpha\ket{0}\otimes\big(\gamma\ket{0}+\delta\ket{1}\big)
    +\beta\ket{1}\otimes\big(\gamma\ket{0}+\delta\ket{1}\big)\\
  &=\alpha\gamma\ket{00}+\alpha\delta\ket{01}+\beta\gamma\ket{10}+\beta\delta\ket{11}.
\end{align*}
Par les composantes, $\begin{pmatrix}\alpha\\ \beta\end{pmatrix}\otimes\begin{pmatrix}\gamma\\ \delta\end{pmatrix}
=\begin{pmatrix}\alpha\gamma\\ \alpha\delta\\ \beta\gamma\\ \beta\delta\end{pmatrix}$ : c'est le même
résultat. La norme est conservée :
\[
  |\alpha\gamma|^2+|\alpha\delta|^2+|\beta\gamma|^2+|\beta\delta|^2
  =\big(|\alpha|^2+|\beta|^2\big)\big(|\gamma|^2+|\delta|^2\big)=1\cdot1=1 .
\]

\paragraph{Exemples.}
\begin{align*}
  \ket{+}\otimes\ket{+}&=\tfrac12\big(\ket{0}+\ket{1}\big)\otimes\big(\ket{0}+\ket{1}\big)
    =\tfrac12\big(\ket{00}+\ket{01}+\ket{10}+\ket{11}\big)
    =\tfrac12\begin{pmatrix}1\\ 1\\ 1\\ 1\end{pmatrix},\\[1ex]
  \ket{+}\otimes\ket{-}&=\tfrac12\big(\ket{0}+\ket{1}\big)\otimes\big(\ket{0}-\ket{1}\big)
    =\tfrac12\big(\ket{00}-\ket{01}+\ket{10}-\ket{11}\big)
    =\tfrac12\begin{pmatrix}1\\ -1\\ 1\\ -1\end{pmatrix},\\[1ex]
  \ket{0}\otimes\ket{+}&=\tfrac{1}{\sqrt2}\big(\ket{00}+\ket{01}\big)
    =\tfrac{1}{\sqrt2}\begin{pmatrix}1\\ 1\\ 0\\ 0\end{pmatrix},\qquad
  \ket{+}\otimes\ket{0}=\tfrac{1}{\sqrt2}\big(\ket{00}+\ket{10}\big)
    =\tfrac{1}{\sqrt2}\begin{pmatrix}1\\ 0\\ 1\\ 0\end{pmatrix}.
\end{align*}
Les deux derniers vecteurs sont différents : \emph{l'ordre des facteurs compte}, le premier
est toujours celui de $A$.

% ---------------------------------------------------------------------
\subsection{États à deux qubits et mesure}\label{sec:mesure2}
% ---------------------------------------------------------------------

Un état quelconque de $\mathcal{H}^4$ se décompose sur la base de calcul :
\[
  \ket{\Psi}=c_{00}\ket{00}+c_{01}\ket{01}+c_{10}\ket{10}+c_{11}\ket{11}
  =\begin{pmatrix}c_{00}\\ c_{01}\\ c_{10}\\ c_{11}\end{pmatrix},
  \qquad \sum_{i,j}|c_{ij}|^2=1 .
\]

\begin{resultat}{Mesure des deux qubits dans la base de calcul}
\[
  P(ij)=|\braket{ij}{\Psi}|^2=|c_{ij}|^2 .
\]
Pour un seul qubit : $P(A{=}0)=P(00)+P(01)$, $P(A{=}1)=P(10)+P(11)$,
$P(B{=}0)=P(00)+P(10)$, $P(B{=}1)=P(01)+P(11)$.
\end{resultat}

En effet $\braket{ij}{\Psi}=\sum_{k,l}c_{kl}\braket{ij}{kl}=\sum_{k,l}c_{kl}\delta_{ik}\delta_{jl}=c_{ij}$.

\paragraph{Exemple 1.} $\ket{\Psi}=\frac{1}{\sqrt2}\ket{00}+\frac{1}{\sqrt2}\ket{01}
=\big(\tfrac{1}{\sqrt2},\tfrac{1}{\sqrt2},0,0\big)^{\mathsf T}$. Il est normalisé :
$\frac12+\frac12+0+0=1$. Les probabilités jointes sont
$P(00)=P(01)=\frac12$ et $P(10)=P(11)=0$. Pour chaque qubit :
\[
  P(A{=}0)=\tfrac12+\tfrac12=1,\quad P(A{=}1)=0+0=0,\quad
  P(B{=}0)=\tfrac12+0=\tfrac12,\quad P(B{=}1)=\tfrac12+0=\tfrac12 .
\]
$A$ donne toujours $0$ ; $B$ donne $0$ ou $1$ au hasard.

\paragraph{Exemple 2 (mesure de $A$ seul).} Soit
$\ket{\Psi}=\frac12\ket{00}+\frac12\ket{01}+\frac{1}{\sqrt2}\ket{11}$. Normalisation :
$\frac14+\frac14+0+\frac12=1$. Probabilités jointes : $P(00)=\frac14$, $P(01)=\frac14$,
$P(10)=0$, $P(11)=\frac12$. Donc
\[
  P(A{=}0)=\tfrac14+\tfrac14=\tfrac12,\quad P(A{=}1)=0+\tfrac12=\tfrac12,\quad
  P(B{=}0)=\tfrac14+0=\tfrac14,\quad P(B{=}1)=\tfrac14+\tfrac12=\tfrac34 .
\]

\begin{remarque}[Effondrement après la mesure de $A$]
Si $A$ donne le résultat $a$, on ne garde que les termes de $\ket{\Psi}$ dont le premier
chiffre vaut $a$, puis on renormalise :
\[
  \ket{\Psi}\ \to\ \frac{c_{a0}\ket{a0}+c_{a1}\ket{a1}}{\sqrt{|c_{a0}|^2+|c_{a1}|^2}},
  \qquad\text{probabilité }|c_{a0}|^2+|c_{a1}|^2 .
\]
\end{remarque}

Pour l'exemple 2 :
\begin{itemize}[nosep]
  \item $A=0$ (probabilité $\frac12$) : on garde $\frac12\ket{00}+\frac12\ket{01}$, de norme
        au carré $\frac12$ ; renormalisé, $\frac{1}{\sqrt2}\big(\ket{00}+\ket{01}\big)
        =\ket{0}\otimes\ket{+}$. Ensuite $B$ vaut $0$ ou $1$ avec la probabilité $\frac12$.
  \item $A=1$ (probabilité $\frac12$) : on garde $\frac{1}{\sqrt2}\ket{11}$, de norme au carré
        $\frac12$ ; renormalisé, $\ket{11}=\ket{1}\otimes\ket{1}$. Ensuite $B$ vaut $1$
        avec certitude.
\end{itemize}
Contrôle : $P(B{=}0)=\frac12\cdot\frac12+\frac12\cdot0=\frac14$, comme calculé plus haut.
Ici la mesure de $A$ \emph{modifie} l'état de $B$ ($\ket{+}$ ou $\ket{1}$) : c'est la
signature d'un état intriqué (section suivante).

% ---------------------------------------------------------------------
\subsection{États séparables et états intriqués}\label{sec:separable}
% ---------------------------------------------------------------------

\begin{definition}{Séparable ou intriqué}
Un état $\ket{\Psi}_{AB}$ est \emph{séparable} s'il s'écrit comme un produit
$\ket{\Psi}_{AB}=\ket{\Psi}_A\otimes\ket{\Psi}_B$. Sinon il est \emph{intriqué} : on ne
peut pas attribuer d'état propre à chaque qubit.
\end{definition}

\paragraph{Exemple séparable (exemple 1 ci-dessus).} Par linéarité du produit tensoriel :
\begin{align*}
  \ket{\Psi}&=\tfrac{1}{\sqrt2}\ket{00}_{AB}+\tfrac{1}{\sqrt2}\ket{01}_{AB}
    =\tfrac{1}{\sqrt2}\Big(\ket{0}_A\otimes\ket{0}_B+\ket{0}_A\otimes\ket{1}_B\Big)\\
  &=\ket{0}_A\otimes\Big(\tfrac{1}{\sqrt2}\big(\ket{0}+\ket{1}\big)\Big)_B
    =\ket{0}_A\otimes\ket{+}_B ,
\end{align*}
donc $\ket{\Psi}_A=\ket{0}$ et $\ket{\Psi}_B=\ket{+}$. Contrôle par les composantes :
$\begin{pmatrix}1\\ 0\end{pmatrix}\otimes\tfrac{1}{\sqrt2}\begin{pmatrix}1\\ 1\end{pmatrix}
=\tfrac{1}{\sqrt2}\begin{pmatrix}1\cdot1\\ 1\cdot1\\ 0\cdot1\\ 0\cdot1\end{pmatrix}
=\tfrac{1}{\sqrt2}\begin{pmatrix}1\\ 1\\ 0\\ 0\end{pmatrix}$.

\subsubsection{Un critère de séparabilité}

Un produit quelconque a pour coefficients $c_{00}=ac$, $c_{01}=ad$, $c_{10}=bc$,
$c_{11}=bd$ (formule de la section~\ref{sec:tenseur}, avec $a,b$ pour $A$ et $c,d$ pour $B$), donc
\[
  c_{00}c_{11}=(ac)(bd)=abcd=(ad)(bc)=c_{01}c_{10}.
\]

\begin{resultat}{Test de séparabilité (deux qubits)}
$\ket{\Psi}=\sum c_{ij}\ket{ij}$ est séparable si et seulement si
\[
  D=c_{00}c_{11}-c_{01}c_{10}=0 .
\]
\end{resultat}

\emph{Démonstration.} ($\Rightarrow$) Calcul ci-dessus. ($\Leftarrow$) On regroupe les
termes selon l'état de $A$ :
\[
  \ket{\Psi}=\ket{0}\otimes\big(c_{00}\ket{0}+c_{01}\ket{1}\big)+\ket{1}\otimes\big(c_{10}\ket{0}+c_{11}\ket{1}\big).
\]
Si $D=0$, les lignes $(c_{00},c_{01})$ et $(c_{10},c_{11})$ sont proportionnelles. Si
$(c_{00},c_{01})\neq(0,0)$, il existe $\lambda$ avec $(c_{10},c_{11})=\lambda(c_{00},c_{01})$
et $\ket{\Psi}=\big(\ket{0}+\lambda\ket{1}\big)\otimes\big(c_{00}\ket{0}+c_{01}\ket{1}\big)$.
Si $(c_{00},c_{01})=(0,0)$, alors $\ket{\Psi}=\ket{1}\otimes\big(c_{10}\ket{0}+c_{11}\ket{1}\big)$.
Dans les deux cas $\ket{\Psi}$ est un produit ; on renormalise chaque facteur.
\hfill$\square$

\subsubsection{Application à plusieurs états}

\begin{center}
\renewcommand{\arraystretch}{1.5}
\begin{tabular}{@{}lccl@{}}
\toprule
\textbf{État} & $(c_{00},c_{01},c_{10},c_{11})$ & $D$ & \textbf{Nature} \\
\midrule
$\ket{00}$ & $(1,0,0,0)$ & $1\cdot0-0\cdot0=0$ & séparable \\
$\ket{+}\otimes\ket{+}$ & $\frac12(1,1,1,1)$ & $\frac14-\frac14=0$ & séparable \\
$\ket{+}\otimes\ket{-}$ & $\frac12(1,-1,1,-1)$ & $-\frac14-(-\frac14)=0$ & séparable \\
$\frac{1}{\sqrt2}(\ket{00}+\ket{01})$ & $\frac{1}{\sqrt2}(1,1,0,0)$ & $0-\frac12\cdot0=0$ & séparable \\
$\ket{\Phi^+}=\frac{1}{\sqrt2}(\ket{00}+\ket{11})$ & $\frac{1}{\sqrt2}(1,0,0,1)$ & $\frac12-0=\frac12$ & intriqué \\
$\ket{\Phi^-}=\frac{1}{\sqrt2}(\ket{00}-\ket{11})$ & $\frac{1}{\sqrt2}(1,0,0,-1)$ & $-\frac12-0=-\frac12$ & intriqué \\
$\ket{\Psi^+}=\frac{1}{\sqrt2}(\ket{01}+\ket{10})$ & $\frac{1}{\sqrt2}(0,1,1,0)$ & $0-\frac12=-\frac12$ & intriqué \\
$\ket{\Psi^-}=\frac{1}{\sqrt2}(\ket{01}-\ket{10})$ & $\frac{1}{\sqrt2}(0,1,-1,0)$ & $0-(-\frac12)=\frac12$ & intriqué \\
$\frac12\ket{00}+\frac12\ket{01}+\frac{1}{\sqrt2}\ket{11}$ & $(\frac12,\frac12,0,\frac{1}{\sqrt2})$ & $\frac{1}{2\sqrt2}-0$ & intriqué \\
\bottomrule
\end{tabular}
\end{center}

Les quatre états de Bell $\ket{\Phi^\pm}$, $\ket{\Psi^\pm}$ sont \emph{orthonormés} :
chacun est de norme $\frac12+\frac12=1$ ; $\ket{\Phi^\pm}$ n'ont des composantes que sur
$\ket{00}$ et $\ket{11}$, $\ket{\Psi^\pm}$ que sur $\ket{01}$ et $\ket{10}$, donc
$\braket{\Phi^\pm}{\Psi^\pm}=0$, et
$\braket{\Phi^+}{\Phi^-}=\frac12(1-1)=0$, $\braket{\Psi^+}{\Psi^-}=\frac12(1-1)=0$.
Ils forment une base de $\mathcal{H}^4$ constituée uniquement d'états intriqués.

\paragraph{Preuve directe que $\ket{\Phi^+}$ est intriqué.} Supposons
$\ket{\Phi^+}=(a\ket{0}+b\ket{1})\otimes(c\ket{0}+d\ket{1})$. En identifiant les
composantes : $ac=\frac{1}{\sqrt2}$, $ad=0$, $bc=0$, $bd=\frac{1}{\sqrt2}$. La condition
$ad=0$ impose $a=0$ (impossible car $ac\neq0$) ou $d=0$ (impossible car $bd\neq0$).
Contradiction : aucun produit ne convient.

\paragraph{Corrélations.} Pour un produit $\ket{\Psi}_A\otimes\ket{\Psi}_B$ avec
$\ket{\Psi}_A=a_0\ket{0}+a_1\ket{1}$ et $\ket{\Psi}_B=b_0\ket{0}+b_1\ket{1}$, on a
$P(ij)=|a_ib_j|^2=|a_i|^2|b_j|^2=P(A{=}i)\,P(B{=}j)$ : les deux qubits sont
\emph{indépendants}. Vérification avec $\ket{+}\otimes\ket{+}$ : $P(00)=\frac14=\frac12\cdot\frac12$.
Pour $\ket{\Phi^+}$ : $P(00)=P(11)=\frac12$ et $P(01)=P(10)=0$, donc
$P(A{=}0)=P(B{=}0)=\frac12$, mais
\[
  P(00)=\tfrac12\neq\tfrac12\cdot\tfrac12=P(A{=}0)\,P(B{=}0) .
\]
Les deux qubits donnent toujours le même résultat, aléatoire : ils sont corrélés.

\paragraph{État d'un qubit pris seul.} On prend la matrice de densité
$\rho=\ket{\Psi}\bra{\Psi}$ ($4\times4$, lignes et colonnes ordonnées $00,01,10,11$) et on
trace sur $B$ : $(\rho_A)_{ij}=\rho_{(i0),(j0)}+\rho_{(i1),(j1)}$.

\emph{État de Bell $\ket{\Phi^+}$.}
\[
  \rho=\frac12\begin{pmatrix}1 & 0 & 0 & 1\\ 0 & 0 & 0 & 0\\ 0 & 0 & 0 & 0\\ 1 & 0 & 0 & 1\end{pmatrix},
  \qquad
  \begin{aligned}
    (\rho_A)_{00}&=\rho_{00,00}+\rho_{01,01}=\tfrac12+0=\tfrac12,\\
    (\rho_A)_{01}&=\rho_{00,10}+\rho_{01,11}=0+0=0,\\
    (\rho_A)_{11}&=\rho_{10,10}+\rho_{11,11}=0+\tfrac12=\tfrac12 .
  \end{aligned}
\]
Donc $\rho_A=\frac12\mathrm{Id}$ : état \emph{mixte} ($\operatorname{Tr}\rho_A^2=\frac12$), au centre
de la boule de Bloch (section~\ref{sec:bloch}).

\emph{État séparable $\ket{0}\otimes\ket{+}$.}
\[
  \rho=\frac12\begin{pmatrix}1 & 1 & 0 & 0\\ 1 & 1 & 0 & 0\\ 0 & 0 & 0 & 0\\ 0 & 0 & 0 & 0\end{pmatrix},
  \qquad
  \begin{aligned}
    (\rho_A)_{00}&=\rho_{00,00}+\rho_{01,01}=\tfrac12+\tfrac12=1,\\
    (\rho_A)_{01}&=\rho_{00,10}+\rho_{01,11}=0+0=0,\\
    (\rho_A)_{11}&=\rho_{10,10}+\rho_{11,11}=0+0=0 .
  \end{aligned}
\]
Donc $\rho_A=\ket{0}\bra{0}$ : état \emph{pur}, celui du facteur $\ket{\Psi}_A$.

Un qubit d'un état séparable garde son propre état pur ; un qubit d'un état de Bell est
complètement aléatoire alors que le couple est parfaitement défini.

% ---------------------------------------------------------------------
\subsection{Opérateurs sur deux qubits}\label{sec:op2}
% ---------------------------------------------------------------------

\begin{definition}{Produit tensoriel de deux matrices}
\[
  A\otimes B=\begin{pmatrix}A_{11} & A_{12}\\ A_{21} & A_{22}\end{pmatrix}\otimes B
  =\begin{pmatrix}A_{11}B & A_{12}B\\ A_{21}B & A_{22}B\end{pmatrix}
  \qquad(\text{matrice }4\times4).
\]
\end{definition}

\begin{resultat}{Action sur un produit}
\[
  (A\otimes B)\big(\ket{u}\otimes\ket{v}\big)=A\ket{u}\otimes B\ket{v}.
\]
\end{resultat}

\emph{Démonstration.} Avec $\ket{u}=(u_1,u_2)^{\mathsf T}$, le vecteur $\ket{u}\otimes\ket{v}$ a
pour blocs $u_k\ket{v}$ ($k=1,2$). Le bloc de lignes $i$ de $(A\otimes B)(\ket{u}\otimes\ket{v})$ est
donc
\[
  \sum_kA_{ik}B\,\big(u_k\ket{v}\big)=\Big(\sum_kA_{ik}u_k\Big)B\ket{v}=(A\ket{u})_i\,B\ket{v},
\]
qui est le bloc $i$ de $A\ket{u}\otimes B\ket{v}$. \hfill$\square$

\subsubsection{Portes agissant sur un seul des deux qubits}

On note $I=\mathrm{Id}$ ($2\times2$). $X\otimes I$ agit sur $A$, $I\otimes X$ agit sur $B$ :
\begin{align*}
  I\otimes X&=\begin{pmatrix}1\cdot X & 0\cdot X\\ 0\cdot X & 1\cdot X\end{pmatrix}
    =\begin{pmatrix}0 & 1 & 0 & 0\\ 1 & 0 & 0 & 0\\ 0 & 0 & 0 & 1\\ 0 & 0 & 1 & 0\end{pmatrix},\\[1ex]
  X\otimes I&=\begin{pmatrix}0\cdot I & 1\cdot I\\ 1\cdot I & 0\cdot I\end{pmatrix}
    =\begin{pmatrix}0 & 0 & 1 & 0\\ 0 & 0 & 0 & 1\\ 1 & 0 & 0 & 0\\ 0 & 1 & 0 & 0\end{pmatrix}.
\end{align*}
Sur $\ket{00}=(1,0,0,0)^{\mathsf T}$ (première colonne de la matrice) :
$(I\otimes X)\ket{00}=(0,1,0,0)^{\mathsf T}=\ket{01}$ et
$(X\otimes I)\ket{00}=(0,0,1,0)^{\mathsf T}=\ket{10}$. Avec la règle :
$(I\otimes X)\ket{00}=I\ket{0}\otimes X\ket{0}=\ket{0}\otimes\ket{1}=\ket{01}$.

Pour la porte de Hadamard, avec $H=\frac{1}{\sqrt2}\begin{pmatrix}1 & 1\\ 1 & -1\end{pmatrix}$ :
\[
  H\otimes I=\frac{1}{\sqrt2}\begin{pmatrix}1\cdot I & 1\cdot I\\ 1\cdot I & -1\cdot I\end{pmatrix}
  =\frac{1}{\sqrt2}\begin{pmatrix}1 & 0 & 1 & 0\\ 0 & 1 & 0 & 1\\ 1 & 0 & -1 & 0\\ 0 & 1 & 0 & -1\end{pmatrix}.
\]
Ses colonnes donnent l'image des quatre états de base :
\begin{align*}
  (H\otimes I)\ket{00}&=\tfrac{1}{\sqrt2}(1,0,1,0)^{\mathsf T}=\ket{+}\otimes\ket{0},&
  (H\otimes I)\ket{01}&=\tfrac{1}{\sqrt2}(0,1,0,1)^{\mathsf T}=\ket{+}\otimes\ket{1},\\
  (H\otimes I)\ket{10}&=\tfrac{1}{\sqrt2}(1,0,-1,0)^{\mathsf T}=\ket{-}\otimes\ket{0},&
  (H\otimes I)\ket{11}&=\tfrac{1}{\sqrt2}(0,1,0,-1)^{\mathsf T}=\ket{-}\otimes\ket{1}.
\end{align*}

\subsubsection{Hadamard sur les deux qubits}

\[
  H\otimes H=\frac12\begin{pmatrix}1\cdot H' & 1\cdot H'\\ 1\cdot H' & -1\cdot H'\end{pmatrix},
  \quad H'=\begin{pmatrix}1 & 1\\ 1 & -1\end{pmatrix},
  \qquad
  H\otimes H=\frac12\begin{pmatrix}1 & 1 & 1 & 1\\ 1 & -1 & 1 & -1\\ 1 & 1 & -1 & -1\\ 1 & -1 & -1 & 1\end{pmatrix}.
\]
La première colonne donne $(H\otimes H)\ket{00}=\frac12(1,1,1,1)^{\mathsf T}=\ket{+}\otimes\ket{+}$
(cohérent avec $H\ket{0}\otimes H\ket{0}$) : les quatre résultats $00,01,10,11$ sont
équiprobables, chacun avec la probabilité $\frac14$. Sur $n$ qubits, $H^{\otimes n}\ket{0\ldots0}$
donne une superposition uniforme de $2^n$ états.

\subsubsection{La porte CNOT et les états de Bell}

\begin{definition}{Porte CNOT (NOT contrôlé)}
Le qubit $A$ est le \emph{contrôle}, $B$ la \emph{cible} : $B$ est inversé si $A=1$.
\[
  \mathrm{CNOT}=\ket{0}\bra{0}\otimes I+\ket{1}\bra{1}\otimes X
  =\begin{pmatrix}I & 0\\ 0 & X\end{pmatrix}
  =\begin{pmatrix}1 & 0 & 0 & 0\\ 0 & 1 & 0 & 0\\ 0 & 0 & 0 & 1\\ 0 & 0 & 1 & 0\end{pmatrix}.
\]
\end{definition}

Sa matrice échange les composantes $10$ et $11$ :
$\mathrm{CNOT}\ket{00}=\ket{00}$, $\mathrm{CNOT}\ket{01}=\ket{01}$,
$\mathrm{CNOT}\ket{10}=\ket{11}$, $\mathrm{CNOT}\ket{11}=\ket{10}$, soit
$\ket{i,j}\to\ket{i,i\oplus j}$ (addition modulo $2$).

On applique $H\otimes I$ puis $\mathrm{CNOT}$ à $\ket{00}$ (figure~\ref{fig:bell}) :
\[
  \ket{00}\ \xrightarrow{\ H\otimes I\ }\ \tfrac{1}{\sqrt2}\begin{pmatrix}1\\ 0\\ 1\\ 0\end{pmatrix}
  =\tfrac{1}{\sqrt2}\big(\ket{00}+\ket{10}\big)
  \ \xrightarrow{\ \mathrm{CNOT}\ }\ \tfrac{1}{\sqrt2}\begin{pmatrix}1\\ 0\\ 0\\ 1\end{pmatrix}
  =\tfrac{1}{\sqrt2}\big(\ket{00}+\ket{11}\big)=\ket{\Phi^+}.
\]

\begin{figure}[!ht]
\centering
\begin{tikzpicture}[>=Stealth, font=\small]
  % fils : B en haut (dernier chiffre du ket), A en bas (premier chiffre)
  \draw (0,1) -- (8.6,1);
  \draw (0,0) -- (8.6,0);
  \node[anchor=east] at (0,1) {$\ket{0}_B$};
  \node[anchor=east] at (0,0) {$\ket{0}_A$};
  % Hadamard sur A
  \node[draw=insarouge, fill=insarouge!8, minimum size=7mm, thick] at (2.2,0) {$H$};
  % CNOT : contrôle sur A (bas), cible sur B (haut)
  \draw[insarouge, thick] (5.4,0) -- (5.4,1.3);
  \fill[insarouge] (5.4,0) circle (2.6pt);
  \draw[insarouge, thick, fill=white] (5.4,1) circle (0.3);
  \draw[insarouge, thick] (5.1,1) -- (5.7,1);
  \draw[insarouge, thick] (5.4,0.7) -- (5.4,1.3);
  % accolade de sortie
  \draw[insagris, thick] (8.75,1.2) -- (8.9,1.2) -- (8.9,-0.2) -- (8.75,-0.2);
  \node[anchor=west] at (8.9,0.5) {$\ket{\Phi^+}_{AB}$};
  % états intermédiaires
  \foreach \x in {0.6,3.8,7.2}
    \draw[dashed, insagris!70] (\x,1.35) -- (\x,-0.95);
  \node[anchor=north, font=\footnotesize] at (0.6,-0.95) {$\ket{00}$};
  \node[anchor=north, font=\footnotesize] at (3.8,-0.95) {$\tfrac{1}{\sqrt2}(\ket{00}+\ket{10})$};
  \node[anchor=north, font=\footnotesize] at (7.2,-0.95) {$\tfrac{1}{\sqrt2}(\ket{00}+\ket{11})$};
\end{tikzpicture}
\caption{Circuit qui produit l'état de Bell $\ket{\Phi^+}$ : une porte de Hadamard sur $A$,
puis un CNOT ($A$ contrôle, $B$ cible). Les états intermédiaires sont indiqués sous les
traits pointillés. Convention : les fils sont dessinés dans l'ordre inverse de l'écriture
du ket $\ket{ij}_{AB}$, si bien que le premier chiffre (celui de $A$) est sur le fil du bas.}
\label{fig:bell}
\end{figure}

Les quatre états de base d'entrée donnent les quatre états de Bell :
\begin{center}
\renewcommand{\arraystretch}{1.6}
\begin{tabular}{@{}cccc@{}}
\toprule
Entrée & après $H\otimes I$ & après CNOT & Sortie \\
\midrule
$\ket{00}$ & $\frac{1}{\sqrt2}(\ket{00}+\ket{10})$ & $\frac{1}{\sqrt2}(\ket{00}+\ket{11})$ & $\ket{\Phi^+}$ \\
$\ket{01}$ & $\frac{1}{\sqrt2}(\ket{01}+\ket{11})$ & $\frac{1}{\sqrt2}(\ket{01}+\ket{10})$ & $\ket{\Psi^+}$ \\
$\ket{10}$ & $\frac{1}{\sqrt2}(\ket{00}-\ket{10})$ & $\frac{1}{\sqrt2}(\ket{00}-\ket{11})$ & $\ket{\Phi^-}$ \\
$\ket{11}$ & $\frac{1}{\sqrt2}(\ket{01}-\ket{11})$ & $\frac{1}{\sqrt2}(\ket{01}-\ket{10})$ & $\ket{\Psi^-}$ \\
\bottomrule
\end{tabular}
\end{center}
Sous forme matricielle, le produit $\mathrm{CNOT}\cdot(H\otimes I)$ échange les lignes $3$ et $4$ de
$H\otimes I$ :
\[
  \mathrm{CNOT}\,(H\otimes I)=\frac{1}{\sqrt2}\begin{pmatrix}1 & 0 & 1 & 0\\ 0 & 1 & 0 & 1\\ 0 & 1 & 0 & -1\\ 1 & 0 & -1 & 0\end{pmatrix},
\]
et ses quatre colonnes sont $\ket{\Phi^+}$, $\ket{\Psi^+}$, $\ket{\Phi^-}$, $\ket{\Psi^-}$.

\begin{remarque}[Porte intriquante]
Les états d'entrée $\ket{ij}$ sont séparables ($D=0$) ; les sorties sont intriquées
($D=\pm\frac12$). Aucune porte agissant séparément sur $A$ et sur $B$ ($U_A\otimes U_B$)
ne peut créer de l'intrication à partir d'un produit, puisque
$(U_A\otimes U_B)(\ket{u}\otimes\ket{v})=U_A\ket{u}\otimes U_B\ket{v}$ reste un produit.
Une porte à deux qubits comme le CNOT est indispensable.
\end{remarque}

% ---------------------------------------------------------------------
\subsection{Récapitulatif du chapitre}
% ---------------------------------------------------------------------

\begin{center}
\renewcommand{\arraystretch}{1.4}
\begin{tabular}{@{}ll@{}}
\toprule
\textbf{Notion} & \textbf{Formule} \\
\midrule
Espace de deux qubits & $\mathcal{H}^2\otimes\mathcal{H}^2=\mathcal{H}^4$, \quad dimension $2^n$ pour $n$ qubits \\
Produit de vecteurs & $(a_1,a_2)^{\mathsf T}\otimes(b_1,b_2)^{\mathsf T}=(a_1b_1,a_1b_2,a_2b_1,a_2b_2)^{\mathsf T}$ \\
Base de calcul & $\ket{ij}_{AB}=\ket{i}_A\otimes\ket{j}_B$, \quad $\braket{ij}{kl}=\delta_{ik}\delta_{jl}$ \\
État général & $\ket{\Psi}=\sum c_{ij}\ket{ij}$, \quad $\sum|c_{ij}|^2=1$, \quad $P(ij)=|c_{ij}|^2$ \\
Séparable & $\ket{\Psi}_A\otimes\ket{\Psi}_B$ \quad $\Leftrightarrow$ \quad $c_{00}c_{11}-c_{01}c_{10}=0$ \\
Intriqué & non séparable, \quad par exemple $\ket{\Phi^+}=\frac{1}{\sqrt2}(\ket{00}+\ket{11})$ \\
Matrices & $A\otimes B=(A_{ij}B)$, \quad $(A\otimes B)(\ket{u}\otimes\ket{v})=A\ket{u}\otimes B\ket{v}$ \\
CNOT & $\ket{i,j}\to\ket{i,i\oplus j}$, \quad $\mathrm{CNOT}\,(H\otimes I)\ket{00}=\ket{\Phi^+}$ \\
\bottomrule
\end{tabular}
\end{center}
.
%  Comme chapitre4.tex et chapitre5.tex, il ouvre une NOUVELLE SECTION
%  (numéro 4). Pour le rattacher à la section précédente, remplacer
%  \section par \subsection et chaque \subsection par \subsubsection.
%  Pas de préambule : environnements (definition, resultat, remarque),
%  macros (\ii, \ee, \dd, \ket, \bra, \braket) et couleurs (insarouge,
%  insagris) sont définis dans main.tex.
% =====================================================================

% Macros de Dirac (déjà définies dans main.tex récent ; ce bloc évite l'erreur
% « Undefined control sequence \ket » si main.tex est une ancienne version).
\providecommand{\ket}[1]{|#1\rangle}
\providecommand{\bra}[1]{\langle #1|}
\providecommand{\braket}[2]{\langle #1|#2\rangle}

% =====================================================================
\section{Systèmes multi-qubits et intrication}\label{sec:multi}
% =====================================================================

Un ordinateur quantique utilise plusieurs qubits. Cette section décrit l'espace des états
de deux qubits (le produit tensoriel), la mesure, la distinction entre états séparables
et états intriqués, puis les opérateurs qui agissent sur deux qubits. Tous les calculs
sont détaillés ; les qubits sont notés $A$ et $B$.

% ---------------------------------------------------------------------
\subsection{L'espace de deux qubits}\label{sec:tenseur}
% ---------------------------------------------------------------------

Chaque qubit vit dans $\mathcal{H}^2=\mathbb{C}^2$. L'espace du système $(A,B)$ est le
\emph{produit tensoriel}
\[
  \mathcal{H}^2\otimes\mathcal{H}^2=\mathcal{H}^4=\mathbb{C}^4 ,
\]
de dimension $2\times2=4$. Pour $n$ qubits la dimension est $2^n$ :

\begin{center}
\renewcommand{\arraystretch}{1.3}
\begin{tabular}{@{}lccccc@{}}
\toprule
Nombre de qubits $n$ & $1$ & $2$ & $3$ & $10$ & $50$ \\
\midrule
Dimension $2^n$ & $2$ & $4$ & $8$ & $1024$ & $\approx1{,}1\times10^{15}$ \\
\bottomrule
\end{tabular}
\end{center}

\begin{definition}{Produit tensoriel de deux vecteurs}
\[
  \begin{pmatrix}a_1\\ a_2\end{pmatrix}\otimes\begin{pmatrix}b_1\\ b_2\end{pmatrix}
  =\begin{pmatrix}a_1\begin{pmatrix}b_1\\ b_2\end{pmatrix}\\[2ex]
                  a_2\begin{pmatrix}b_1\\ b_2\end{pmatrix}\end{pmatrix}
  =\begin{pmatrix}a_1b_1\\ a_1b_2\\ a_2b_1\\ a_2b_2\end{pmatrix}.
\]
Chaque composante du premier vecteur multiplie le second vecteur en entier. On note
$\ket{ij}_{AB}=\ket{i}_A\otimes\ket{j}_B$ : le premier chiffre est l'état de $A$, le second
celui de $B$.
\end{definition}

\subsubsection{Les quatre états de base}

On développe chacun des quatre produits avec $\ket{0}=\begin{pmatrix}1\\ 0\end{pmatrix}$
et $\ket{1}=\begin{pmatrix}0\\ 1\end{pmatrix}$ :
\begin{align*}
  \ket{0}\otimes\ket{0}&=\begin{pmatrix}1\\ 0\end{pmatrix}\otimes\begin{pmatrix}1\\ 0\end{pmatrix}
    =\begin{pmatrix}1\cdot\begin{pmatrix}1\\ 0\end{pmatrix}\\[1.5ex]0\cdot\begin{pmatrix}1\\ 0\end{pmatrix}\end{pmatrix}
    =\begin{pmatrix}1\\ 0\\ 0\\ 0\end{pmatrix}=\ket{00},\\[1.5ex]
  \ket{0}\otimes\ket{1}&=\begin{pmatrix}1\\ 0\end{pmatrix}\otimes\begin{pmatrix}0\\ 1\end{pmatrix}
    =\begin{pmatrix}1\cdot\begin{pmatrix}0\\ 1\end{pmatrix}\\[1.5ex]0\cdot\begin{pmatrix}0\\ 1\end{pmatrix}\end{pmatrix}
    =\begin{pmatrix}0\\ 1\\ 0\\ 0\end{pmatrix}=\ket{01},\\[1.5ex]
  \ket{1}\otimes\ket{0}&=\begin{pmatrix}0\\ 1\end{pmatrix}\otimes\begin{pmatrix}1\\ 0\end{pmatrix}
    =\begin{pmatrix}0\cdot\begin{pmatrix}1\\ 0\end{pmatrix}\\[1.5ex]1\cdot\begin{pmatrix}1\\ 0\end{pmatrix}\end{pmatrix}
    =\begin{pmatrix}0\\ 0\\ 1\\ 0\end{pmatrix}=\ket{10},\\[1.5ex]
  \ket{1}\otimes\ket{1}&=\begin{pmatrix}0\\ 1\end{pmatrix}\otimes\begin{pmatrix}0\\ 1\end{pmatrix}
    =\begin{pmatrix}0\cdot\begin{pmatrix}0\\ 1\end{pmatrix}\\[1.5ex]1\cdot\begin{pmatrix}0\\ 1\end{pmatrix}\end{pmatrix}
    =\begin{pmatrix}0\\ 0\\ 0\\ 1\end{pmatrix}=\ket{11}.
\end{align*}
Le vecteur $\ket{ij}$ a son $1$ en position $2i+j$ (en comptant à partir de $0$) : c'est
la lecture de $ij$ en binaire. Les quatre vecteurs forment la base de calcul de
$\mathcal{H}^4$.

\subsubsection{Produit scalaire et orthonormalité}

\begin{resultat}{Produit scalaire d'un produit tensoriel}
\[
  \big(\bra{a}\otimes\bra{b}\big)\big(\ket{c}\otimes\ket{d}\big)=\braket{a}{c}\,\braket{b}{d}.
\]
\end{resultat}

\emph{Démonstration.} Les composantes de $\ket{a}\otimes\ket{b}$ sont $a_ib_j$, celles de
$\ket{c}\otimes\ket{d}$ sont $c_id_j$. Le produit scalaire est
$\sum_{i,j}(a_ib_j)^*c_id_j=\big(\sum_ia_i^*c_i\big)\big(\sum_jb_j^*d_j\big)=\braket{a}{c}\braket{b}{d}$.
\hfill$\square$

On en déduit $\braket{ij}{kl}=\braket{i}{k}\braket{j}{l}$. Les seize produits scalaires
des états de base sont :
\begin{center}
\renewcommand{\arraystretch}{1.3}
\begin{tabular}{@{}c|cccc@{}}
 & $\ket{00}$ & $\ket{01}$ & $\ket{10}$ & $\ket{11}$ \\
\hline
$\bra{00}$ & $1\cdot1=1$ & $1\cdot0=0$ & $0\cdot1=0$ & $0\cdot0=0$ \\
$\bra{01}$ & $1\cdot0=0$ & $1\cdot1=1$ & $0\cdot0=0$ & $0\cdot1=0$ \\
$\bra{10}$ & $0\cdot1=0$ & $0\cdot0=0$ & $1\cdot1=1$ & $1\cdot0=0$ \\
$\bra{11}$ & $0\cdot0=0$ & $0\cdot1=0$ & $1\cdot0=0$ & $1\cdot1=1$ \\
\end{tabular}
\end{center}
(chaque case est $\braket{i}{k}\cdot\braket{j}{l}$). On lit $\braket{ij}{kl}=\delta_{ik}\delta_{jl}$ :
la base est \emph{orthonormée}.

\subsubsection{Produit tensoriel d'états quelconques}

Le produit tensoriel est linéaire en chaque facteur :
$(\lambda\ket{a}+\mu\ket{a'})\otimes\ket{b}=\lambda\ket{a}\otimes\ket{b}+\mu\ket{a'}\otimes\ket{b}$,
et de même pour le second facteur. Pour $\alpha\ket{0}+\beta\ket{1}$ et
$\gamma\ket{0}+\delta\ket{1}$ :
\begin{align*}
  \big(\alpha\ket{0}+\beta\ket{1}\big)\otimes\big(\gamma\ket{0}+\delta\ket{1}\big)
  &=\alpha\ket{0}\otimes\big(\gamma\ket{0}+\delta\ket{1}\big)
    +\beta\ket{1}\otimes\big(\gamma\ket{0}+\delta\ket{1}\big)\\
  &=\alpha\gamma\ket{00}+\alpha\delta\ket{01}+\beta\gamma\ket{10}+\beta\delta\ket{11}.
\end{align*}
Par les composantes, $\begin{pmatrix}\alpha\\ \beta\end{pmatrix}\otimes\begin{pmatrix}\gamma\\ \delta\end{pmatrix}
=\begin{pmatrix}\alpha\gamma\\ \alpha\delta\\ \beta\gamma\\ \beta\delta\end{pmatrix}$ : c'est le même
résultat. La norme est conservée :
\[
  |\alpha\gamma|^2+|\alpha\delta|^2+|\beta\gamma|^2+|\beta\delta|^2
  =\big(|\alpha|^2+|\beta|^2\big)\big(|\gamma|^2+|\delta|^2\big)=1\cdot1=1 .
\]

\paragraph{Exemples.}
\begin{align*}
  \ket{+}\otimes\ket{+}&=\tfrac12\big(\ket{0}+\ket{1}\big)\otimes\big(\ket{0}+\ket{1}\big)
    =\tfrac12\big(\ket{00}+\ket{01}+\ket{10}+\ket{11}\big)
    =\tfrac12\begin{pmatrix}1\\ 1\\ 1\\ 1\end{pmatrix},\\[1ex]
  \ket{+}\otimes\ket{-}&=\tfrac12\big(\ket{0}+\ket{1}\big)\otimes\big(\ket{0}-\ket{1}\big)
    =\tfrac12\big(\ket{00}-\ket{01}+\ket{10}-\ket{11}\big)
    =\tfrac12\begin{pmatrix}1\\ -1\\ 1\\ -1\end{pmatrix},\\[1ex]
  \ket{0}\otimes\ket{+}&=\tfrac{1}{\sqrt2}\big(\ket{00}+\ket{01}\big)
    =\tfrac{1}{\sqrt2}\begin{pmatrix}1\\ 1\\ 0\\ 0\end{pmatrix},\qquad
  \ket{+}\otimes\ket{0}=\tfrac{1}{\sqrt2}\big(\ket{00}+\ket{10}\big)
    =\tfrac{1}{\sqrt2}\begin{pmatrix}1\\ 0\\ 1\\ 0\end{pmatrix}.
\end{align*}
Les deux derniers vecteurs sont différents : \emph{l'ordre des facteurs compte}, le premier
est toujours celui de $A$.

% ---------------------------------------------------------------------
\subsection{États à deux qubits et mesure}\label{sec:mesure2}
% ---------------------------------------------------------------------

Un état quelconque de $\mathcal{H}^4$ se décompose sur la base de calcul :
\[
  \ket{\Psi}=c_{00}\ket{00}+c_{01}\ket{01}+c_{10}\ket{10}+c_{11}\ket{11}
  =\begin{pmatrix}c_{00}\\ c_{01}\\ c_{10}\\ c_{11}\end{pmatrix},
  \qquad \sum_{i,j}|c_{ij}|^2=1 .
\]

\begin{resultat}{Mesure des deux qubits dans la base de calcul}
\[
  P(ij)=|\braket{ij}{\Psi}|^2=|c_{ij}|^2 .
\]
Pour un seul qubit : $P(A{=}0)=P(00)+P(01)$, $P(A{=}1)=P(10)+P(11)$,
$P(B{=}0)=P(00)+P(10)$, $P(B{=}1)=P(01)+P(11)$.
\end{resultat}

En effet $\braket{ij}{\Psi}=\sum_{k,l}c_{kl}\braket{ij}{kl}=\sum_{k,l}c_{kl}\delta_{ik}\delta_{jl}=c_{ij}$.

\paragraph{Exemple 1.} $\ket{\Psi}=\frac{1}{\sqrt2}\ket{00}+\frac{1}{\sqrt2}\ket{01}
=\big(\tfrac{1}{\sqrt2},\tfrac{1}{\sqrt2},0,0\big)^{\mathsf T}$. Il est normalisé :
$\frac12+\frac12+0+0=1$. Les probabilités jointes sont
$P(00)=P(01)=\frac12$ et $P(10)=P(11)=0$. Pour chaque qubit :
\[
  P(A{=}0)=\tfrac12+\tfrac12=1,\quad P(A{=}1)=0+0=0,\quad
  P(B{=}0)=\tfrac12+0=\tfrac12,\quad P(B{=}1)=\tfrac12+0=\tfrac12 .
\]
$A$ donne toujours $0$ ; $B$ donne $0$ ou $1$ au hasard.

\paragraph{Exemple 2 (mesure de $A$ seul).} Soit
$\ket{\Psi}=\frac12\ket{00}+\frac12\ket{01}+\frac{1}{\sqrt2}\ket{11}$. Normalisation :
$\frac14+\frac14+0+\frac12=1$. Probabilités jointes : $P(00)=\frac14$, $P(01)=\frac14$,
$P(10)=0$, $P(11)=\frac12$. Donc
\[
  P(A{=}0)=\tfrac14+\tfrac14=\tfrac12,\quad P(A{=}1)=0+\tfrac12=\tfrac12,\quad
  P(B{=}0)=\tfrac14+0=\tfrac14,\quad P(B{=}1)=\tfrac14+\tfrac12=\tfrac34 .
\]

\begin{remarque}[Effondrement après la mesure de $A$]
Si $A$ donne le résultat $a$, on ne garde que les termes de $\ket{\Psi}$ dont le premier
chiffre vaut $a$, puis on renormalise :
\[
  \ket{\Psi}\ \to\ \frac{c_{a0}\ket{a0}+c_{a1}\ket{a1}}{\sqrt{|c_{a0}|^2+|c_{a1}|^2}},
  \qquad\text{probabilité }|c_{a0}|^2+|c_{a1}|^2 .
\]
\end{remarque}

Pour l'exemple 2 :
\begin{itemize}[nosep]
  \item $A=0$ (probabilité $\frac12$) : on garde $\frac12\ket{00}+\frac12\ket{01}$, de norme
        au carré $\frac12$ ; renormalisé, $\frac{1}{\sqrt2}\big(\ket{00}+\ket{01}\big)
        =\ket{0}\otimes\ket{+}$. Ensuite $B$ vaut $0$ ou $1$ avec la probabilité $\frac12$.
  \item $A=1$ (probabilité $\frac12$) : on garde $\frac{1}{\sqrt2}\ket{11}$, de norme au carré
        $\frac12$ ; renormalisé, $\ket{11}=\ket{1}\otimes\ket{1}$. Ensuite $B$ vaut $1$
        avec certitude.
\end{itemize}
Contrôle : $P(B{=}0)=\frac12\cdot\frac12+\frac12\cdot0=\frac14$, comme calculé plus haut.
Ici la mesure de $A$ \emph{modifie} l'état de $B$ ($\ket{+}$ ou $\ket{1}$) : c'est la
signature d'un état intriqué (section suivante).

% ---------------------------------------------------------------------
\subsection{États séparables et états intriqués}\label{sec:separable}
% ---------------------------------------------------------------------

\begin{definition}{Séparable ou intriqué}
Un état $\ket{\Psi}_{AB}$ est \emph{séparable} s'il s'écrit comme un produit
$\ket{\Psi}_{AB}=\ket{\Psi}_A\otimes\ket{\Psi}_B$. Sinon il est \emph{intriqué} : on ne
peut pas attribuer d'état propre à chaque qubit.
\end{definition}

\paragraph{Exemple séparable (exemple 1 ci-dessus).} Par linéarité du produit tensoriel :
\begin{align*}
  \ket{\Psi}&=\tfrac{1}{\sqrt2}\ket{00}_{AB}+\tfrac{1}{\sqrt2}\ket{01}_{AB}
    =\tfrac{1}{\sqrt2}\Big(\ket{0}_A\otimes\ket{0}_B+\ket{0}_A\otimes\ket{1}_B\Big)\\
  &=\ket{0}_A\otimes\Big(\tfrac{1}{\sqrt2}\big(\ket{0}+\ket{1}\big)\Big)_B
    =\ket{0}_A\otimes\ket{+}_B ,
\end{align*}
donc $\ket{\Psi}_A=\ket{0}$ et $\ket{\Psi}_B=\ket{+}$. Contrôle par les composantes :
$\begin{pmatrix}1\\ 0\end{pmatrix}\otimes\tfrac{1}{\sqrt2}\begin{pmatrix}1\\ 1\end{pmatrix}
=\tfrac{1}{\sqrt2}\begin{pmatrix}1\cdot1\\ 1\cdot1\\ 0\cdot1\\ 0\cdot1\end{pmatrix}
=\tfrac{1}{\sqrt2}\begin{pmatrix}1\\ 1\\ 0\\ 0\end{pmatrix}$.

\subsubsection{Un critère de séparabilité}

Un produit quelconque a pour coefficients $c_{00}=ac$, $c_{01}=ad$, $c_{10}=bc$,
$c_{11}=bd$ (formule de la section~\ref{sec:tenseur}, avec $a,b$ pour $A$ et $c,d$ pour $B$), donc
\[
  c_{00}c_{11}=(ac)(bd)=abcd=(ad)(bc)=c_{01}c_{10}.
\]

\begin{resultat}{Test de séparabilité (deux qubits)}
$\ket{\Psi}=\sum c_{ij}\ket{ij}$ est séparable si et seulement si
\[
  D=c_{00}c_{11}-c_{01}c_{10}=0 .
\]
\end{resultat}

\emph{Démonstration.} ($\Rightarrow$) Calcul ci-dessus. ($\Leftarrow$) On regroupe les
termes selon l'état de $A$ :
\[
  \ket{\Psi}=\ket{0}\otimes\big(c_{00}\ket{0}+c_{01}\ket{1}\big)+\ket{1}\otimes\big(c_{10}\ket{0}+c_{11}\ket{1}\big).
\]
Si $D=0$, les lignes $(c_{00},c_{01})$ et $(c_{10},c_{11})$ sont proportionnelles. Si
$(c_{00},c_{01})\neq(0,0)$, il existe $\lambda$ avec $(c_{10},c_{11})=\lambda(c_{00},c_{01})$
et $\ket{\Psi}=\big(\ket{0}+\lambda\ket{1}\big)\otimes\big(c_{00}\ket{0}+c_{01}\ket{1}\big)$.
Si $(c_{00},c_{01})=(0,0)$, alors $\ket{\Psi}=\ket{1}\otimes\big(c_{10}\ket{0}+c_{11}\ket{1}\big)$.
Dans les deux cas $\ket{\Psi}$ est un produit ; on renormalise chaque facteur.
\hfill$\square$

\subsubsection{Application à plusieurs états}

\begin{center}
\renewcommand{\arraystretch}{1.5}
\begin{tabular}{@{}lccl@{}}
\toprule
\textbf{État} & $(c_{00},c_{01},c_{10},c_{11})$ & $D$ & \textbf{Nature} \\
\midrule
$\ket{00}$ & $(1,0,0,0)$ & $1\cdot0-0\cdot0=0$ & séparable \\
$\ket{+}\otimes\ket{+}$ & $\frac12(1,1,1,1)$ & $\frac14-\frac14=0$ & séparable \\
$\ket{+}\otimes\ket{-}$ & $\frac12(1,-1,1,-1)$ & $-\frac14-(-\frac14)=0$ & séparable \\
$\frac{1}{\sqrt2}(\ket{00}+\ket{01})$ & $\frac{1}{\sqrt2}(1,1,0,0)$ & $0-\frac12\cdot0=0$ & séparable \\
$\ket{\Phi^+}=\frac{1}{\sqrt2}(\ket{00}+\ket{11})$ & $\frac{1}{\sqrt2}(1,0,0,1)$ & $\frac12-0=\frac12$ & intriqué \\
$\ket{\Phi^-}=\frac{1}{\sqrt2}(\ket{00}-\ket{11})$ & $\frac{1}{\sqrt2}(1,0,0,-1)$ & $-\frac12-0=-\frac12$ & intriqué \\
$\ket{\Psi^+}=\frac{1}{\sqrt2}(\ket{01}+\ket{10})$ & $\frac{1}{\sqrt2}(0,1,1,0)$ & $0-\frac12=-\frac12$ & intriqué \\
$\ket{\Psi^-}=\frac{1}{\sqrt2}(\ket{01}-\ket{10})$ & $\frac{1}{\sqrt2}(0,1,-1,0)$ & $0-(-\frac12)=\frac12$ & intriqué \\
$\frac12\ket{00}+\frac12\ket{01}+\frac{1}{\sqrt2}\ket{11}$ & $(\frac12,\frac12,0,\frac{1}{\sqrt2})$ & $\frac{1}{2\sqrt2}-0$ & intriqué \\
\bottomrule
\end{tabular}
\end{center}

Les quatre états de Bell $\ket{\Phi^\pm}$, $\ket{\Psi^\pm}$ sont \emph{orthonormés} :
chacun est de norme $\frac12+\frac12=1$ ; $\ket{\Phi^\pm}$ n'ont des composantes que sur
$\ket{00}$ et $\ket{11}$, $\ket{\Psi^\pm}$ que sur $\ket{01}$ et $\ket{10}$, donc
$\braket{\Phi^\pm}{\Psi^\pm}=0$, et
$\braket{\Phi^+}{\Phi^-}=\frac12(1-1)=0$, $\braket{\Psi^+}{\Psi^-}=\frac12(1-1)=0$.
Ils forment une base de $\mathcal{H}^4$ constituée uniquement d'états intriqués.

\paragraph{Preuve directe que $\ket{\Phi^+}$ est intriqué.} Supposons
$\ket{\Phi^+}=(a\ket{0}+b\ket{1})\otimes(c\ket{0}+d\ket{1})$. En identifiant les
composantes : $ac=\frac{1}{\sqrt2}$, $ad=0$, $bc=0$, $bd=\frac{1}{\sqrt2}$. La condition
$ad=0$ impose $a=0$ (impossible car $ac\neq0$) ou $d=0$ (impossible car $bd\neq0$).
Contradiction : aucun produit ne convient.

\paragraph{Corrélations.} Pour un produit $\ket{\Psi}_A\otimes\ket{\Psi}_B$ avec
$\ket{\Psi}_A=a_0\ket{0}+a_1\ket{1}$ et $\ket{\Psi}_B=b_0\ket{0}+b_1\ket{1}$, on a
$P(ij)=|a_ib_j|^2=|a_i|^2|b_j|^2=P(A{=}i)\,P(B{=}j)$ : les deux qubits sont
\emph{indépendants}. Vérification avec $\ket{+}\otimes\ket{+}$ : $P(00)=\frac14=\frac12\cdot\frac12$.
Pour $\ket{\Phi^+}$ : $P(00)=P(11)=\frac12$ et $P(01)=P(10)=0$, donc
$P(A{=}0)=P(B{=}0)=\frac12$, mais
\[
  P(00)=\tfrac12\neq\tfrac12\cdot\tfrac12=P(A{=}0)\,P(B{=}0) .
\]
Les deux qubits donnent toujours le même résultat, aléatoire : ils sont corrélés.

\paragraph{État d'un qubit pris seul.} On prend la matrice de densité
$\rho=\ket{\Psi}\bra{\Psi}$ ($4\times4$, lignes et colonnes ordonnées $00,01,10,11$) et on
trace sur $B$ : $(\rho_A)_{ij}=\rho_{(i0),(j0)}+\rho_{(i1),(j1)}$.

\emph{État de Bell $\ket{\Phi^+}$.}
\[
  \rho=\frac12\begin{pmatrix}1 & 0 & 0 & 1\\ 0 & 0 & 0 & 0\\ 0 & 0 & 0 & 0\\ 1 & 0 & 0 & 1\end{pmatrix},
  \qquad
  \begin{aligned}
    (\rho_A)_{00}&=\rho_{00,00}+\rho_{01,01}=\tfrac12+0=\tfrac12,\\
    (\rho_A)_{01}&=\rho_{00,10}+\rho_{01,11}=0+0=0,\\
    (\rho_A)_{11}&=\rho_{10,10}+\rho_{11,11}=0+\tfrac12=\tfrac12 .
  \end{aligned}
\]
Donc $\rho_A=\frac12\mathrm{Id}$ : état \emph{mixte} ($\operatorname{Tr}\rho_A^2=\frac12$), au centre
de la boule de Bloch (section~\ref{sec:bloch}).

\emph{État séparable $\ket{0}\otimes\ket{+}$.}
\[
  \rho=\frac12\begin{pmatrix}1 & 1 & 0 & 0\\ 1 & 1 & 0 & 0\\ 0 & 0 & 0 & 0\\ 0 & 0 & 0 & 0\end{pmatrix},
  \qquad
  \begin{aligned}
    (\rho_A)_{00}&=\rho_{00,00}+\rho_{01,01}=\tfrac12+\tfrac12=1,\\
    (\rho_A)_{01}&=\rho_{00,10}+\rho_{01,11}=0+0=0,\\
    (\rho_A)_{11}&=\rho_{10,10}+\rho_{11,11}=0+0=0 .
  \end{aligned}
\]
Donc $\rho_A=\ket{0}\bra{0}$ : état \emph{pur}, celui du facteur $\ket{\Psi}_A$.

Un qubit d'un état séparable garde son propre état pur ; un qubit d'un état de Bell est
complètement aléatoire alors que le couple est parfaitement défini.

% ---------------------------------------------------------------------
\subsection{Opérateurs sur deux qubits}\label{sec:op2}
% ---------------------------------------------------------------------

\begin{definition}{Produit tensoriel de deux matrices}
\[
  A\otimes B=\begin{pmatrix}A_{11} & A_{12}\\ A_{21} & A_{22}\end{pmatrix}\otimes B
  =\begin{pmatrix}A_{11}B & A_{12}B\\ A_{21}B & A_{22}B\end{pmatrix}
  \qquad(\text{matrice }4\times4).
\]
\end{definition}

\begin{resultat}{Action sur un produit}
\[
  (A\otimes B)\big(\ket{u}\otimes\ket{v}\big)=A\ket{u}\otimes B\ket{v}.
\]
\end{resultat}

\emph{Démonstration.} Avec $\ket{u}=(u_1,u_2)^{\mathsf T}$, le vecteur $\ket{u}\otimes\ket{v}$ a
pour blocs $u_k\ket{v}$ ($k=1,2$). Le bloc de lignes $i$ de $(A\otimes B)(\ket{u}\otimes\ket{v})$ est
donc
\[
  \sum_kA_{ik}B\,\big(u_k\ket{v}\big)=\Big(\sum_kA_{ik}u_k\Big)B\ket{v}=(A\ket{u})_i\,B\ket{v},
\]
qui est le bloc $i$ de $A\ket{u}\otimes B\ket{v}$. \hfill$\square$

\subsubsection{Portes agissant sur un seul des deux qubits}

On note $I=\mathrm{Id}$ ($2\times2$). $X\otimes I$ agit sur $A$, $I\otimes X$ agit sur $B$ :
\begin{align*}
  I\otimes X&=\begin{pmatrix}1\cdot X & 0\cdot X\\ 0\cdot X & 1\cdot X\end{pmatrix}
    =\begin{pmatrix}0 & 1 & 0 & 0\\ 1 & 0 & 0 & 0\\ 0 & 0 & 0 & 1\\ 0 & 0 & 1 & 0\end{pmatrix},\\[1ex]
  X\otimes I&=\begin{pmatrix}0\cdot I & 1\cdot I\\ 1\cdot I & 0\cdot I\end{pmatrix}
    =\begin{pmatrix}0 & 0 & 1 & 0\\ 0 & 0 & 0 & 1\\ 1 & 0 & 0 & 0\\ 0 & 1 & 0 & 0\end{pmatrix}.
\end{align*}
Sur $\ket{00}=(1,0,0,0)^{\mathsf T}$ (première colonne de la matrice) :
$(I\otimes X)\ket{00}=(0,1,0,0)^{\mathsf T}=\ket{01}$ et
$(X\otimes I)\ket{00}=(0,0,1,0)^{\mathsf T}=\ket{10}$. Avec la règle :
$(I\otimes X)\ket{00}=I\ket{0}\otimes X\ket{0}=\ket{0}\otimes\ket{1}=\ket{01}$.

Pour la porte de Hadamard, avec $H=\frac{1}{\sqrt2}\begin{pmatrix}1 & 1\\ 1 & -1\end{pmatrix}$ :
\[
  H\otimes I=\frac{1}{\sqrt2}\begin{pmatrix}1\cdot I & 1\cdot I\\ 1\cdot I & -1\cdot I\end{pmatrix}
  =\frac{1}{\sqrt2}\begin{pmatrix}1 & 0 & 1 & 0\\ 0 & 1 & 0 & 1\\ 1 & 0 & -1 & 0\\ 0 & 1 & 0 & -1\end{pmatrix}.
\]
Ses colonnes donnent l'image des quatre états de base :
\begin{align*}
  (H\otimes I)\ket{00}&=\tfrac{1}{\sqrt2}(1,0,1,0)^{\mathsf T}=\ket{+}\otimes\ket{0},&
  (H\otimes I)\ket{01}&=\tfrac{1}{\sqrt2}(0,1,0,1)^{\mathsf T}=\ket{+}\otimes\ket{1},\\
  (H\otimes I)\ket{10}&=\tfrac{1}{\sqrt2}(1,0,-1,0)^{\mathsf T}=\ket{-}\otimes\ket{0},&
  (H\otimes I)\ket{11}&=\tfrac{1}{\sqrt2}(0,1,0,-1)^{\mathsf T}=\ket{-}\otimes\ket{1}.
\end{align*}

\subsubsection{Hadamard sur les deux qubits}

\[
  H\otimes H=\frac12\begin{pmatrix}1\cdot H' & 1\cdot H'\\ 1\cdot H' & -1\cdot H'\end{pmatrix},
  \quad H'=\begin{pmatrix}1 & 1\\ 1 & -1\end{pmatrix},
  \qquad
  H\otimes H=\frac12\begin{pmatrix}1 & 1 & 1 & 1\\ 1 & -1 & 1 & -1\\ 1 & 1 & -1 & -1\\ 1 & -1 & -1 & 1\end{pmatrix}.
\]
La première colonne donne $(H\otimes H)\ket{00}=\frac12(1,1,1,1)^{\mathsf T}=\ket{+}\otimes\ket{+}$
(cohérent avec $H\ket{0}\otimes H\ket{0}$) : les quatre résultats $00,01,10,11$ sont
équiprobables, chacun avec la probabilité $\frac14$. Sur $n$ qubits, $H^{\otimes n}\ket{0\ldots0}$
donne une superposition uniforme de $2^n$ états.

\subsubsection{La porte CNOT et les états de Bell}

\begin{definition}{Porte CNOT (NOT contrôlé)}
Le qubit $A$ est le \emph{contrôle}, $B$ la \emph{cible} : $B$ est inversé si $A=1$.
\[
  \mathrm{CNOT}=\ket{0}\bra{0}\otimes I+\ket{1}\bra{1}\otimes X
  =\begin{pmatrix}I & 0\\ 0 & X\end{pmatrix}
  =\begin{pmatrix}1 & 0 & 0 & 0\\ 0 & 1 & 0 & 0\\ 0 & 0 & 0 & 1\\ 0 & 0 & 1 & 0\end{pmatrix}.
\]
\end{definition}

Sa matrice échange les composantes $10$ et $11$ :
$\mathrm{CNOT}\ket{00}=\ket{00}$, $\mathrm{CNOT}\ket{01}=\ket{01}$,
$\mathrm{CNOT}\ket{10}=\ket{11}$, $\mathrm{CNOT}\ket{11}=\ket{10}$, soit
$\ket{i,j}\to\ket{i,i\oplus j}$ (addition modulo $2$).

On applique $H\otimes I$ puis $\mathrm{CNOT}$ à $\ket{00}$ (figure~\ref{fig:bell}) :
\[
  \ket{00}\ \xrightarrow{\ H\otimes I\ }\ \tfrac{1}{\sqrt2}\begin{pmatrix}1\\ 0\\ 1\\ 0\end{pmatrix}
  =\tfrac{1}{\sqrt2}\big(\ket{00}+\ket{10}\big)
  \ \xrightarrow{\ \mathrm{CNOT}\ }\ \tfrac{1}{\sqrt2}\begin{pmatrix}1\\ 0\\ 0\\ 1\end{pmatrix}
  =\tfrac{1}{\sqrt2}\big(\ket{00}+\ket{11}\big)=\ket{\Phi^+}.
\]

\begin{figure}[!ht]
\centering
\begin{tikzpicture}[>=Stealth, font=\small]
  % fils : B en haut (dernier chiffre du ket), A en bas (premier chiffre)
  \draw (0,1) -- (8.6,1);
  \draw (0,0) -- (8.6,0);
  \node[anchor=east] at (0,1) {$\ket{0}_B$};
  \node[anchor=east] at (0,0) {$\ket{0}_A$};
  % Hadamard sur A
  \node[draw=insarouge, fill=insarouge!8, minimum size=7mm, thick] at (2.2,0) {$H$};
  % CNOT : contrôle sur A (bas), cible sur B (haut)
  \draw[insarouge, thick] (5.4,0) -- (5.4,1.3);
  \fill[insarouge] (5.4,0) circle (2.6pt);
  \draw[insarouge, thick, fill=white] (5.4,1) circle (0.3);
  \draw[insarouge, thick] (5.1,1) -- (5.7,1);
  \draw[insarouge, thick] (5.4,0.7) -- (5.4,1.3);
  % accolade de sortie
  \draw[insagris, thick] (8.75,1.2) -- (8.9,1.2) -- (8.9,-0.2) -- (8.75,-0.2);
  \node[anchor=west] at (8.9,0.5) {$\ket{\Phi^+}_{AB}$};
  % états intermédiaires
  \foreach \x in {0.6,3.8,7.2}
    \draw[dashed, insagris!70] (\x,1.35) -- (\x,-0.95);
  \node[anchor=north, font=\footnotesize] at (0.6,-0.95) {$\ket{00}$};
  \node[anchor=north, font=\footnotesize] at (3.8,-0.95) {$\tfrac{1}{\sqrt2}(\ket{00}+\ket{10})$};
  \node[anchor=north, font=\footnotesize] at (7.2,-0.95) {$\tfrac{1}{\sqrt2}(\ket{00}+\ket{11})$};
\end{tikzpicture}
\caption{Circuit qui produit l'état de Bell $\ket{\Phi^+}$ : une porte de Hadamard sur $A$,
puis un CNOT ($A$ contrôle, $B$ cible). Les états intermédiaires sont indiqués sous les
traits pointillés. Convention : les fils sont dessinés dans l'ordre inverse de l'écriture
du ket $\ket{ij}_{AB}$, si bien que le premier chiffre (celui de $A$) est sur le fil du bas.}
\label{fig:bell}
\end{figure}

Les quatre états de base d'entrée donnent les quatre états de Bell :
\begin{center}
\renewcommand{\arraystretch}{1.6}
\begin{tabular}{@{}cccc@{}}
\toprule
Entrée & après $H\otimes I$ & après CNOT & Sortie \\
\midrule
$\ket{00}$ & $\frac{1}{\sqrt2}(\ket{00}+\ket{10})$ & $\frac{1}{\sqrt2}(\ket{00}+\ket{11})$ & $\ket{\Phi^+}$ \\
$\ket{01}$ & $\frac{1}{\sqrt2}(\ket{01}+\ket{11})$ & $\frac{1}{\sqrt2}(\ket{01}+\ket{10})$ & $\ket{\Psi^+}$ \\
$\ket{10}$ & $\frac{1}{\sqrt2}(\ket{00}-\ket{10})$ & $\frac{1}{\sqrt2}(\ket{00}-\ket{11})$ & $\ket{\Phi^-}$ \\
$\ket{11}$ & $\frac{1}{\sqrt2}(\ket{01}-\ket{11})$ & $\frac{1}{\sqrt2}(\ket{01}-\ket{10})$ & $\ket{\Psi^-}$ \\
\bottomrule
\end{tabular}
\end{center}
Sous forme matricielle, le produit $\mathrm{CNOT}\cdot(H\otimes I)$ échange les lignes $3$ et $4$ de
$H\otimes I$ :
\[
  \mathrm{CNOT}\,(H\otimes I)=\frac{1}{\sqrt2}\begin{pmatrix}1 & 0 & 1 & 0\\ 0 & 1 & 0 & 1\\ 0 & 1 & 0 & -1\\ 1 & 0 & -1 & 0\end{pmatrix},
\]
et ses quatre colonnes sont $\ket{\Phi^+}$, $\ket{\Psi^+}$, $\ket{\Phi^-}$, $\ket{\Psi^-}$.

\begin{remarque}[Porte intriquante]
Les états d'entrée $\ket{ij}$ sont séparables ($D=0$) ; les sorties sont intriquées
($D=\pm\frac12$). Aucune porte agissant séparément sur $A$ et sur $B$ ($U_A\otimes U_B$)
ne peut créer de l'intrication à partir d'un produit, puisque
$(U_A\otimes U_B)(\ket{u}\otimes\ket{v})=U_A\ket{u}\otimes U_B\ket{v}$ reste un produit.
Une porte à deux qubits comme le CNOT est indispensable.
\end{remarque}

% ---------------------------------------------------------------------
\subsection{Récapitulatif du chapitre}
% ---------------------------------------------------------------------

\begin{center}
\renewcommand{\arraystretch}{1.4}
\begin{tabular}{@{}ll@{}}
\toprule
\textbf{Notion} & \textbf{Formule} \\
\midrule
Espace de deux qubits & $\mathcal{H}^2\otimes\mathcal{H}^2=\mathcal{H}^4$, \quad dimension $2^n$ pour $n$ qubits \\
Produit de vecteurs & $(a_1,a_2)^{\mathsf T}\otimes(b_1,b_2)^{\mathsf T}=(a_1b_1,a_1b_2,a_2b_1,a_2b_2)^{\mathsf T}$ \\
Base de calcul & $\ket{ij}_{AB}=\ket{i}_A\otimes\ket{j}_B$, \quad $\braket{ij}{kl}=\delta_{ik}\delta_{jl}$ \\
État général & $\ket{\Psi}=\sum c_{ij}\ket{ij}$, \quad $\sum|c_{ij}|^2=1$, \quad $P(ij)=|c_{ij}|^2$ \\
Séparable & $\ket{\Psi}_A\otimes\ket{\Psi}_B$ \quad $\Leftrightarrow$ \quad $c_{00}c_{11}-c_{01}c_{10}=0$ \\
Intriqué & non séparable, \quad par exemple $\ket{\Phi^+}=\frac{1}{\sqrt2}(\ket{00}+\ket{11})$ \\
Matrices & $A\otimes B=(A_{ij}B)$, \quad $(A\otimes B)(\ket{u}\otimes\ket{v})=A\ket{u}\otimes B\ket{v}$ \\
CNOT & $\ket{i,j}\to\ket{i,i\oplus j}$, \quad $\mathrm{CNOT}\,(H\otimes I)\ket{00}=\ket{\Phi^+}$ \\
\bottomrule
\end{tabular}
\end{center}
.
%  Comme chapitre4.tex et chapitre5.tex, il ouvre une NOUVELLE SECTION
%  (numéro 4). Pour le rattacher à la section précédente, remplacer
%  \section par \subsection et chaque \subsection par \subsubsection.
%  Pas de préambule : environnements (definition, resultat, remarque),
%  macros (\ii, \ee, \dd, \ket, \bra, \braket) et couleurs (insarouge,
%  insagris) sont définis dans main.tex.
% =====================================================================

% Macros de Dirac (déjà définies dans main.tex récent ; ce bloc évite l'erreur
% « Undefined control sequence \ket » si main.tex est une ancienne version).
\providecommand{\ket}[1]{|#1\rangle}
\providecommand{\bra}[1]{\langle #1|}
\providecommand{\braket}[2]{\langle #1|#2\rangle}

% =====================================================================
\section{Systèmes multi-qubits et intrication}\label{sec:multi}
% =====================================================================

Un ordinateur quantique utilise plusieurs qubits. Cette section décrit l'espace des états
de deux qubits (le produit tensoriel), la mesure, la distinction entre états séparables
et états intriqués, puis les opérateurs qui agissent sur deux qubits. Tous les calculs
sont détaillés ; les qubits sont notés $A$ et $B$.

% ---------------------------------------------------------------------
\subsection{L'espace de deux qubits}\label{sec:tenseur}
% ---------------------------------------------------------------------

Chaque qubit vit dans $\mathcal{H}^2=\mathbb{C}^2$. L'espace du système $(A,B)$ est le
\emph{produit tensoriel}
\[
  \mathcal{H}^2\otimes\mathcal{H}^2=\mathcal{H}^4=\mathbb{C}^4 ,
\]
de dimension $2\times2=4$. Pour $n$ qubits la dimension est $2^n$ :

\begin{center}
\renewcommand{\arraystretch}{1.3}
\begin{tabular}{@{}lccccc@{}}
\toprule
Nombre de qubits $n$ & $1$ & $2$ & $3$ & $10$ & $50$ \\
\midrule
Dimension $2^n$ & $2$ & $4$ & $8$ & $1024$ & $\approx1{,}1\times10^{15}$ \\
\bottomrule
\end{tabular}
\end{center}

\begin{definition}{Produit tensoriel de deux vecteurs}
\[
  \begin{pmatrix}a_1\\ a_2\end{pmatrix}\otimes\begin{pmatrix}b_1\\ b_2\end{pmatrix}
  =\begin{pmatrix}a_1\begin{pmatrix}b_1\\ b_2\end{pmatrix}\\[2ex]
                  a_2\begin{pmatrix}b_1\\ b_2\end{pmatrix}\end{pmatrix}
  =\begin{pmatrix}a_1b_1\\ a_1b_2\\ a_2b_1\\ a_2b_2\end{pmatrix}.
\]
Chaque composante du premier vecteur multiplie le second vecteur en entier. On note
$\ket{ij}_{AB}=\ket{i}_A\otimes\ket{j}_B$ : le premier chiffre est l'état de $A$, le second
celui de $B$.
\end{definition}

\subsubsection{Les quatre états de base}

On développe chacun des quatre produits avec $\ket{0}=\begin{pmatrix}1\\ 0\end{pmatrix}$
et $\ket{1}=\begin{pmatrix}0\\ 1\end{pmatrix}$ :
\begin{align*}
  \ket{0}\otimes\ket{0}&=\begin{pmatrix}1\\ 0\end{pmatrix}\otimes\begin{pmatrix}1\\ 0\end{pmatrix}
    =\begin{pmatrix}1\cdot\begin{pmatrix}1\\ 0\end{pmatrix}\\[1.5ex]0\cdot\begin{pmatrix}1\\ 0\end{pmatrix}\end{pmatrix}
    =\begin{pmatrix}1\\ 0\\ 0\\ 0\end{pmatrix}=\ket{00},\\[1.5ex]
  \ket{0}\otimes\ket{1}&=\begin{pmatrix}1\\ 0\end{pmatrix}\otimes\begin{pmatrix}0\\ 1\end{pmatrix}
    =\begin{pmatrix}1\cdot\begin{pmatrix}0\\ 1\end{pmatrix}\\[1.5ex]0\cdot\begin{pmatrix}0\\ 1\end{pmatrix}\end{pmatrix}
    =\begin{pmatrix}0\\ 1\\ 0\\ 0\end{pmatrix}=\ket{01},\\[1.5ex]
  \ket{1}\otimes\ket{0}&=\begin{pmatrix}0\\ 1\end{pmatrix}\otimes\begin{pmatrix}1\\ 0\end{pmatrix}
    =\begin{pmatrix}0\cdot\begin{pmatrix}1\\ 0\end{pmatrix}\\[1.5ex]1\cdot\begin{pmatrix}1\\ 0\end{pmatrix}\end{pmatrix}
    =\begin{pmatrix}0\\ 0\\ 1\\ 0\end{pmatrix}=\ket{10},\\[1.5ex]
  \ket{1}\otimes\ket{1}&=\begin{pmatrix}0\\ 1\end{pmatrix}\otimes\begin{pmatrix}0\\ 1\end{pmatrix}
    =\begin{pmatrix}0\cdot\begin{pmatrix}0\\ 1\end{pmatrix}\\[1.5ex]1\cdot\begin{pmatrix}0\\ 1\end{pmatrix}\end{pmatrix}
    =\begin{pmatrix}0\\ 0\\ 0\\ 1\end{pmatrix}=\ket{11}.
\end{align*}
Le vecteur $\ket{ij}$ a son $1$ en position $2i+j$ (en comptant à partir de $0$) : c'est
la lecture de $ij$ en binaire. Les quatre vecteurs forment la base de calcul de
$\mathcal{H}^4$.

\subsubsection{Produit scalaire et orthonormalité}

\begin{resultat}{Produit scalaire d'un produit tensoriel}
\[
  \big(\bra{a}\otimes\bra{b}\big)\big(\ket{c}\otimes\ket{d}\big)=\braket{a}{c}\,\braket{b}{d}.
\]
\end{resultat}

\emph{Démonstration.} Les composantes de $\ket{a}\otimes\ket{b}$ sont $a_ib_j$, celles de
$\ket{c}\otimes\ket{d}$ sont $c_id_j$. Le produit scalaire est
$\sum_{i,j}(a_ib_j)^*c_id_j=\big(\sum_ia_i^*c_i\big)\big(\sum_jb_j^*d_j\big)=\braket{a}{c}\braket{b}{d}$.
\hfill$\square$

On en déduit $\braket{ij}{kl}=\braket{i}{k}\braket{j}{l}$. Les seize produits scalaires
des états de base sont :
\begin{center}
\renewcommand{\arraystretch}{1.3}
\begin{tabular}{@{}c|cccc@{}}
 & $\ket{00}$ & $\ket{01}$ & $\ket{10}$ & $\ket{11}$ \\
\hline
$\bra{00}$ & $1\cdot1=1$ & $1\cdot0=0$ & $0\cdot1=0$ & $0\cdot0=0$ \\
$\bra{01}$ & $1\cdot0=0$ & $1\cdot1=1$ & $0\cdot0=0$ & $0\cdot1=0$ \\
$\bra{10}$ & $0\cdot1=0$ & $0\cdot0=0$ & $1\cdot1=1$ & $1\cdot0=0$ \\
$\bra{11}$ & $0\cdot0=0$ & $0\cdot1=0$ & $1\cdot0=0$ & $1\cdot1=1$ \\
\end{tabular}
\end{center}
(chaque case est $\braket{i}{k}\cdot\braket{j}{l}$). On lit $\braket{ij}{kl}=\delta_{ik}\delta_{jl}$ :
la base est \emph{orthonormée}.

\subsubsection{Produit tensoriel d'états quelconques}

Le produit tensoriel est linéaire en chaque facteur :
$(\lambda\ket{a}+\mu\ket{a'})\otimes\ket{b}=\lambda\ket{a}\otimes\ket{b}+\mu\ket{a'}\otimes\ket{b}$,
et de même pour le second facteur. Pour $\alpha\ket{0}+\beta\ket{1}$ et
$\gamma\ket{0}+\delta\ket{1}$ :
\begin{align*}
  \big(\alpha\ket{0}+\beta\ket{1}\big)\otimes\big(\gamma\ket{0}+\delta\ket{1}\big)
  &=\alpha\ket{0}\otimes\big(\gamma\ket{0}+\delta\ket{1}\big)
    +\beta\ket{1}\otimes\big(\gamma\ket{0}+\delta\ket{1}\big)\\
  &=\alpha\gamma\ket{00}+\alpha\delta\ket{01}+\beta\gamma\ket{10}+\beta\delta\ket{11}.
\end{align*}
Par les composantes, $\begin{pmatrix}\alpha\\ \beta\end{pmatrix}\otimes\begin{pmatrix}\gamma\\ \delta\end{pmatrix}
=\begin{pmatrix}\alpha\gamma\\ \alpha\delta\\ \beta\gamma\\ \beta\delta\end{pmatrix}$ : c'est le même
résultat. La norme est conservée :
\[
  |\alpha\gamma|^2+|\alpha\delta|^2+|\beta\gamma|^2+|\beta\delta|^2
  =\big(|\alpha|^2+|\beta|^2\big)\big(|\gamma|^2+|\delta|^2\big)=1\cdot1=1 .
\]

\paragraph{Exemples.}
\begin{align*}
  \ket{+}\otimes\ket{+}&=\tfrac12\big(\ket{0}+\ket{1}\big)\otimes\big(\ket{0}+\ket{1}\big)
    =\tfrac12\big(\ket{00}+\ket{01}+\ket{10}+\ket{11}\big)
    =\tfrac12\begin{pmatrix}1\\ 1\\ 1\\ 1\end{pmatrix},\\[1ex]
  \ket{+}\otimes\ket{-}&=\tfrac12\big(\ket{0}+\ket{1}\big)\otimes\big(\ket{0}-\ket{1}\big)
    =\tfrac12\big(\ket{00}-\ket{01}+\ket{10}-\ket{11}\big)
    =\tfrac12\begin{pmatrix}1\\ -1\\ 1\\ -1\end{pmatrix},\\[1ex]
  \ket{0}\otimes\ket{+}&=\tfrac{1}{\sqrt2}\big(\ket{00}+\ket{01}\big)
    =\tfrac{1}{\sqrt2}\begin{pmatrix}1\\ 1\\ 0\\ 0\end{pmatrix},\qquad
  \ket{+}\otimes\ket{0}=\tfrac{1}{\sqrt2}\big(\ket{00}+\ket{10}\big)
    =\tfrac{1}{\sqrt2}\begin{pmatrix}1\\ 0\\ 1\\ 0\end{pmatrix}.
\end{align*}
Les deux derniers vecteurs sont différents : \emph{l'ordre des facteurs compte}, le premier
est toujours celui de $A$.

% ---------------------------------------------------------------------
\subsection{États à deux qubits et mesure}\label{sec:mesure2}
% ---------------------------------------------------------------------

Un état quelconque de $\mathcal{H}^4$ se décompose sur la base de calcul :
\[
  \ket{\Psi}=c_{00}\ket{00}+c_{01}\ket{01}+c_{10}\ket{10}+c_{11}\ket{11}
  =\begin{pmatrix}c_{00}\\ c_{01}\\ c_{10}\\ c_{11}\end{pmatrix},
  \qquad \sum_{i,j}|c_{ij}|^2=1 .
\]

\begin{resultat}{Mesure des deux qubits dans la base de calcul}
\[
  P(ij)=|\braket{ij}{\Psi}|^2=|c_{ij}|^2 .
\]
Pour un seul qubit : $P(A{=}0)=P(00)+P(01)$, $P(A{=}1)=P(10)+P(11)$,
$P(B{=}0)=P(00)+P(10)$, $P(B{=}1)=P(01)+P(11)$.
\end{resultat}

En effet $\braket{ij}{\Psi}=\sum_{k,l}c_{kl}\braket{ij}{kl}=\sum_{k,l}c_{kl}\delta_{ik}\delta_{jl}=c_{ij}$.

\paragraph{Exemple 1.} $\ket{\Psi}=\frac{1}{\sqrt2}\ket{00}+\frac{1}{\sqrt2}\ket{01}
=\big(\tfrac{1}{\sqrt2},\tfrac{1}{\sqrt2},0,0\big)^{\mathsf T}$. Il est normalisé :
$\frac12+\frac12+0+0=1$. Les probabilités jointes sont
$P(00)=P(01)=\frac12$ et $P(10)=P(11)=0$. Pour chaque qubit :
\[
  P(A{=}0)=\tfrac12+\tfrac12=1,\quad P(A{=}1)=0+0=0,\quad
  P(B{=}0)=\tfrac12+0=\tfrac12,\quad P(B{=}1)=\tfrac12+0=\tfrac12 .
\]
$A$ donne toujours $0$ ; $B$ donne $0$ ou $1$ au hasard.

\paragraph{Exemple 2 (mesure de $A$ seul).} Soit
$\ket{\Psi}=\frac12\ket{00}+\frac12\ket{01}+\frac{1}{\sqrt2}\ket{11}$. Normalisation :
$\frac14+\frac14+0+\frac12=1$. Probabilités jointes : $P(00)=\frac14$, $P(01)=\frac14$,
$P(10)=0$, $P(11)=\frac12$. Donc
\[
  P(A{=}0)=\tfrac14+\tfrac14=\tfrac12,\quad P(A{=}1)=0+\tfrac12=\tfrac12,\quad
  P(B{=}0)=\tfrac14+0=\tfrac14,\quad P(B{=}1)=\tfrac14+\tfrac12=\tfrac34 .
\]

\begin{remarque}[Effondrement après la mesure de $A$]
Si $A$ donne le résultat $a$, on ne garde que les termes de $\ket{\Psi}$ dont le premier
chiffre vaut $a$, puis on renormalise :
\[
  \ket{\Psi}\ \to\ \frac{c_{a0}\ket{a0}+c_{a1}\ket{a1}}{\sqrt{|c_{a0}|^2+|c_{a1}|^2}},
  \qquad\text{probabilité }|c_{a0}|^2+|c_{a1}|^2 .
\]
\end{remarque}

Pour l'exemple 2 :
\begin{itemize}[nosep]
  \item $A=0$ (probabilité $\frac12$) : on garde $\frac12\ket{00}+\frac12\ket{01}$, de norme
        au carré $\frac12$ ; renormalisé, $\frac{1}{\sqrt2}\big(\ket{00}+\ket{01}\big)
        =\ket{0}\otimes\ket{+}$. Ensuite $B$ vaut $0$ ou $1$ avec la probabilité $\frac12$.
  \item $A=1$ (probabilité $\frac12$) : on garde $\frac{1}{\sqrt2}\ket{11}$, de norme au carré
        $\frac12$ ; renormalisé, $\ket{11}=\ket{1}\otimes\ket{1}$. Ensuite $B$ vaut $1$
        avec certitude.
\end{itemize}
Contrôle : $P(B{=}0)=\frac12\cdot\frac12+\frac12\cdot0=\frac14$, comme calculé plus haut.
Ici la mesure de $A$ \emph{modifie} l'état de $B$ ($\ket{+}$ ou $\ket{1}$) : c'est la
signature d'un état intriqué (section suivante).

% ---------------------------------------------------------------------
\subsection{États séparables et états intriqués}\label{sec:separable}
% ---------------------------------------------------------------------

\begin{definition}{Séparable ou intriqué}
Un état $\ket{\Psi}_{AB}$ est \emph{séparable} s'il s'écrit comme un produit
$\ket{\Psi}_{AB}=\ket{\Psi}_A\otimes\ket{\Psi}_B$. Sinon il est \emph{intriqué} : on ne
peut pas attribuer d'état propre à chaque qubit.
\end{definition}

\paragraph{Exemple séparable (exemple 1 ci-dessus).} Par linéarité du produit tensoriel :
\begin{align*}
  \ket{\Psi}&=\tfrac{1}{\sqrt2}\ket{00}_{AB}+\tfrac{1}{\sqrt2}\ket{01}_{AB}
    =\tfrac{1}{\sqrt2}\Big(\ket{0}_A\otimes\ket{0}_B+\ket{0}_A\otimes\ket{1}_B\Big)\\
  &=\ket{0}_A\otimes\Big(\tfrac{1}{\sqrt2}\big(\ket{0}+\ket{1}\big)\Big)_B
    =\ket{0}_A\otimes\ket{+}_B ,
\end{align*}
donc $\ket{\Psi}_A=\ket{0}$ et $\ket{\Psi}_B=\ket{+}$. Contrôle par les composantes :
$\begin{pmatrix}1\\ 0\end{pmatrix}\otimes\tfrac{1}{\sqrt2}\begin{pmatrix}1\\ 1\end{pmatrix}
=\tfrac{1}{\sqrt2}\begin{pmatrix}1\cdot1\\ 1\cdot1\\ 0\cdot1\\ 0\cdot1\end{pmatrix}
=\tfrac{1}{\sqrt2}\begin{pmatrix}1\\ 1\\ 0\\ 0\end{pmatrix}$.

\subsubsection{Un critère de séparabilité}

Un produit quelconque a pour coefficients $c_{00}=ac$, $c_{01}=ad$, $c_{10}=bc$,
$c_{11}=bd$ (formule de la section~\ref{sec:tenseur}, avec $a,b$ pour $A$ et $c,d$ pour $B$), donc
\[
  c_{00}c_{11}=(ac)(bd)=abcd=(ad)(bc)=c_{01}c_{10}.
\]

\begin{resultat}{Test de séparabilité (deux qubits)}
$\ket{\Psi}=\sum c_{ij}\ket{ij}$ est séparable si et seulement si
\[
  D=c_{00}c_{11}-c_{01}c_{10}=0 .
\]
\end{resultat}

\emph{Démonstration.} ($\Rightarrow$) Calcul ci-dessus. ($\Leftarrow$) On regroupe les
termes selon l'état de $A$ :
\[
  \ket{\Psi}=\ket{0}\otimes\big(c_{00}\ket{0}+c_{01}\ket{1}\big)+\ket{1}\otimes\big(c_{10}\ket{0}+c_{11}\ket{1}\big).
\]
Si $D=0$, les lignes $(c_{00},c_{01})$ et $(c_{10},c_{11})$ sont proportionnelles. Si
$(c_{00},c_{01})\neq(0,0)$, il existe $\lambda$ avec $(c_{10},c_{11})=\lambda(c_{00},c_{01})$
et $\ket{\Psi}=\big(\ket{0}+\lambda\ket{1}\big)\otimes\big(c_{00}\ket{0}+c_{01}\ket{1}\big)$.
Si $(c_{00},c_{01})=(0,0)$, alors $\ket{\Psi}=\ket{1}\otimes\big(c_{10}\ket{0}+c_{11}\ket{1}\big)$.
Dans les deux cas $\ket{\Psi}$ est un produit ; on renormalise chaque facteur.
\hfill$\square$

\subsubsection{Application à plusieurs états}

\begin{center}
\renewcommand{\arraystretch}{1.5}
\begin{tabular}{@{}lccl@{}}
\toprule
\textbf{État} & $(c_{00},c_{01},c_{10},c_{11})$ & $D$ & \textbf{Nature} \\
\midrule
$\ket{00}$ & $(1,0,0,0)$ & $1\cdot0-0\cdot0=0$ & séparable \\
$\ket{+}\otimes\ket{+}$ & $\frac12(1,1,1,1)$ & $\frac14-\frac14=0$ & séparable \\
$\ket{+}\otimes\ket{-}$ & $\frac12(1,-1,1,-1)$ & $-\frac14-(-\frac14)=0$ & séparable \\
$\frac{1}{\sqrt2}(\ket{00}+\ket{01})$ & $\frac{1}{\sqrt2}(1,1,0,0)$ & $0-\frac12\cdot0=0$ & séparable \\
$\ket{\Phi^+}=\frac{1}{\sqrt2}(\ket{00}+\ket{11})$ & $\frac{1}{\sqrt2}(1,0,0,1)$ & $\frac12-0=\frac12$ & intriqué \\
$\ket{\Phi^-}=\frac{1}{\sqrt2}(\ket{00}-\ket{11})$ & $\frac{1}{\sqrt2}(1,0,0,-1)$ & $-\frac12-0=-\frac12$ & intriqué \\
$\ket{\Psi^+}=\frac{1}{\sqrt2}(\ket{01}+\ket{10})$ & $\frac{1}{\sqrt2}(0,1,1,0)$ & $0-\frac12=-\frac12$ & intriqué \\
$\ket{\Psi^-}=\frac{1}{\sqrt2}(\ket{01}-\ket{10})$ & $\frac{1}{\sqrt2}(0,1,-1,0)$ & $0-(-\frac12)=\frac12$ & intriqué \\
$\frac12\ket{00}+\frac12\ket{01}+\frac{1}{\sqrt2}\ket{11}$ & $(\frac12,\frac12,0,\frac{1}{\sqrt2})$ & $\frac{1}{2\sqrt2}-0$ & intriqué \\
\bottomrule
\end{tabular}
\end{center}

Les quatre états de Bell $\ket{\Phi^\pm}$, $\ket{\Psi^\pm}$ sont \emph{orthonormés} :
chacun est de norme $\frac12+\frac12=1$ ; $\ket{\Phi^\pm}$ n'ont des composantes que sur
$\ket{00}$ et $\ket{11}$, $\ket{\Psi^\pm}$ que sur $\ket{01}$ et $\ket{10}$, donc
$\braket{\Phi^\pm}{\Psi^\pm}=0$, et
$\braket{\Phi^+}{\Phi^-}=\frac12(1-1)=0$, $\braket{\Psi^+}{\Psi^-}=\frac12(1-1)=0$.
Ils forment une base de $\mathcal{H}^4$ constituée uniquement d'états intriqués.

\paragraph{Preuve directe que $\ket{\Phi^+}$ est intriqué.} Supposons
$\ket{\Phi^+}=(a\ket{0}+b\ket{1})\otimes(c\ket{0}+d\ket{1})$. En identifiant les
composantes : $ac=\frac{1}{\sqrt2}$, $ad=0$, $bc=0$, $bd=\frac{1}{\sqrt2}$. La condition
$ad=0$ impose $a=0$ (impossible car $ac\neq0$) ou $d=0$ (impossible car $bd\neq0$).
Contradiction : aucun produit ne convient.

\paragraph{Corrélations.} Pour un produit $\ket{\Psi}_A\otimes\ket{\Psi}_B$ avec
$\ket{\Psi}_A=a_0\ket{0}+a_1\ket{1}$ et $\ket{\Psi}_B=b_0\ket{0}+b_1\ket{1}$, on a
$P(ij)=|a_ib_j|^2=|a_i|^2|b_j|^2=P(A{=}i)\,P(B{=}j)$ : les deux qubits sont
\emph{indépendants}. Vérification avec $\ket{+}\otimes\ket{+}$ : $P(00)=\frac14=\frac12\cdot\frac12$.
Pour $\ket{\Phi^+}$ : $P(00)=P(11)=\frac12$ et $P(01)=P(10)=0$, donc
$P(A{=}0)=P(B{=}0)=\frac12$, mais
\[
  P(00)=\tfrac12\neq\tfrac12\cdot\tfrac12=P(A{=}0)\,P(B{=}0) .
\]
Les deux qubits donnent toujours le même résultat, aléatoire : ils sont corrélés.

\paragraph{État d'un qubit pris seul.} On prend la matrice de densité
$\rho=\ket{\Psi}\bra{\Psi}$ ($4\times4$, lignes et colonnes ordonnées $00,01,10,11$) et on
trace sur $B$ : $(\rho_A)_{ij}=\rho_{(i0),(j0)}+\rho_{(i1),(j1)}$.

\emph{État de Bell $\ket{\Phi^+}$.}
\[
  \rho=\frac12\begin{pmatrix}1 & 0 & 0 & 1\\ 0 & 0 & 0 & 0\\ 0 & 0 & 0 & 0\\ 1 & 0 & 0 & 1\end{pmatrix},
  \qquad
  \begin{aligned}
    (\rho_A)_{00}&=\rho_{00,00}+\rho_{01,01}=\tfrac12+0=\tfrac12,\\
    (\rho_A)_{01}&=\rho_{00,10}+\rho_{01,11}=0+0=0,\\
    (\rho_A)_{11}&=\rho_{10,10}+\rho_{11,11}=0+\tfrac12=\tfrac12 .
  \end{aligned}
\]
Donc $\rho_A=\frac12\mathrm{Id}$ : état \emph{mixte} ($\operatorname{Tr}\rho_A^2=\frac12$), au centre
de la boule de Bloch (section~\ref{sec:bloch}).

\emph{État séparable $\ket{0}\otimes\ket{+}$.}
\[
  \rho=\frac12\begin{pmatrix}1 & 1 & 0 & 0\\ 1 & 1 & 0 & 0\\ 0 & 0 & 0 & 0\\ 0 & 0 & 0 & 0\end{pmatrix},
  \qquad
  \begin{aligned}
    (\rho_A)_{00}&=\rho_{00,00}+\rho_{01,01}=\tfrac12+\tfrac12=1,\\
    (\rho_A)_{01}&=\rho_{00,10}+\rho_{01,11}=0+0=0,\\
    (\rho_A)_{11}&=\rho_{10,10}+\rho_{11,11}=0+0=0 .
  \end{aligned}
\]
Donc $\rho_A=\ket{0}\bra{0}$ : état \emph{pur}, celui du facteur $\ket{\Psi}_A$.

Un qubit d'un état séparable garde son propre état pur ; un qubit d'un état de Bell est
complètement aléatoire alors que le couple est parfaitement défini.

% ---------------------------------------------------------------------
\subsection{Opérateurs sur deux qubits}\label{sec:op2}
% ---------------------------------------------------------------------

\begin{definition}{Produit tensoriel de deux matrices}
\[
  A\otimes B=\begin{pmatrix}A_{11} & A_{12}\\ A_{21} & A_{22}\end{pmatrix}\otimes B
  =\begin{pmatrix}A_{11}B & A_{12}B\\ A_{21}B & A_{22}B\end{pmatrix}
  \qquad(\text{matrice }4\times4).
\]
\end{definition}

\begin{resultat}{Action sur un produit}
\[
  (A\otimes B)\big(\ket{u}\otimes\ket{v}\big)=A\ket{u}\otimes B\ket{v}.
\]
\end{resultat}

\emph{Démonstration.} Avec $\ket{u}=(u_1,u_2)^{\mathsf T}$, le vecteur $\ket{u}\otimes\ket{v}$ a
pour blocs $u_k\ket{v}$ ($k=1,2$). Le bloc de lignes $i$ de $(A\otimes B)(\ket{u}\otimes\ket{v})$ est
donc
\[
  \sum_kA_{ik}B\,\big(u_k\ket{v}\big)=\Big(\sum_kA_{ik}u_k\Big)B\ket{v}=(A\ket{u})_i\,B\ket{v},
\]
qui est le bloc $i$ de $A\ket{u}\otimes B\ket{v}$. \hfill$\square$

\subsubsection{Portes agissant sur un seul des deux qubits}

On note $I=\mathrm{Id}$ ($2\times2$). $X\otimes I$ agit sur $A$, $I\otimes X$ agit sur $B$ :
\begin{align*}
  I\otimes X&=\begin{pmatrix}1\cdot X & 0\cdot X\\ 0\cdot X & 1\cdot X\end{pmatrix}
    =\begin{pmatrix}0 & 1 & 0 & 0\\ 1 & 0 & 0 & 0\\ 0 & 0 & 0 & 1\\ 0 & 0 & 1 & 0\end{pmatrix},\\[1ex]
  X\otimes I&=\begin{pmatrix}0\cdot I & 1\cdot I\\ 1\cdot I & 0\cdot I\end{pmatrix}
    =\begin{pmatrix}0 & 0 & 1 & 0\\ 0 & 0 & 0 & 1\\ 1 & 0 & 0 & 0\\ 0 & 1 & 0 & 0\end{pmatrix}.
\end{align*}
Sur $\ket{00}=(1,0,0,0)^{\mathsf T}$ (première colonne de la matrice) :
$(I\otimes X)\ket{00}=(0,1,0,0)^{\mathsf T}=\ket{01}$ et
$(X\otimes I)\ket{00}=(0,0,1,0)^{\mathsf T}=\ket{10}$. Avec la règle :
$(I\otimes X)\ket{00}=I\ket{0}\otimes X\ket{0}=\ket{0}\otimes\ket{1}=\ket{01}$.

Pour la porte de Hadamard, avec $H=\frac{1}{\sqrt2}\begin{pmatrix}1 & 1\\ 1 & -1\end{pmatrix}$ :
\[
  H\otimes I=\frac{1}{\sqrt2}\begin{pmatrix}1\cdot I & 1\cdot I\\ 1\cdot I & -1\cdot I\end{pmatrix}
  =\frac{1}{\sqrt2}\begin{pmatrix}1 & 0 & 1 & 0\\ 0 & 1 & 0 & 1\\ 1 & 0 & -1 & 0\\ 0 & 1 & 0 & -1\end{pmatrix}.
\]
Ses colonnes donnent l'image des quatre états de base :
\begin{align*}
  (H\otimes I)\ket{00}&=\tfrac{1}{\sqrt2}(1,0,1,0)^{\mathsf T}=\ket{+}\otimes\ket{0},&
  (H\otimes I)\ket{01}&=\tfrac{1}{\sqrt2}(0,1,0,1)^{\mathsf T}=\ket{+}\otimes\ket{1},\\
  (H\otimes I)\ket{10}&=\tfrac{1}{\sqrt2}(1,0,-1,0)^{\mathsf T}=\ket{-}\otimes\ket{0},&
  (H\otimes I)\ket{11}&=\tfrac{1}{\sqrt2}(0,1,0,-1)^{\mathsf T}=\ket{-}\otimes\ket{1}.
\end{align*}

\subsubsection{Hadamard sur les deux qubits}

\[
  H\otimes H=\frac12\begin{pmatrix}1\cdot H' & 1\cdot H'\\ 1\cdot H' & -1\cdot H'\end{pmatrix},
  \quad H'=\begin{pmatrix}1 & 1\\ 1 & -1\end{pmatrix},
  \qquad
  H\otimes H=\frac12\begin{pmatrix}1 & 1 & 1 & 1\\ 1 & -1 & 1 & -1\\ 1 & 1 & -1 & -1\\ 1 & -1 & -1 & 1\end{pmatrix}.
\]
La première colonne donne $(H\otimes H)\ket{00}=\frac12(1,1,1,1)^{\mathsf T}=\ket{+}\otimes\ket{+}$
(cohérent avec $H\ket{0}\otimes H\ket{0}$) : les quatre résultats $00,01,10,11$ sont
équiprobables, chacun avec la probabilité $\frac14$. Sur $n$ qubits, $H^{\otimes n}\ket{0\ldots0}$
donne une superposition uniforme de $2^n$ états.

\subsubsection{La porte CNOT et les états de Bell}

\begin{definition}{Porte CNOT (NOT contrôlé)}
Le qubit $A$ est le \emph{contrôle}, $B$ la \emph{cible} : $B$ est inversé si $A=1$.
\[
  \mathrm{CNOT}=\ket{0}\bra{0}\otimes I+\ket{1}\bra{1}\otimes X
  =\begin{pmatrix}I & 0\\ 0 & X\end{pmatrix}
  =\begin{pmatrix}1 & 0 & 0 & 0\\ 0 & 1 & 0 & 0\\ 0 & 0 & 0 & 1\\ 0 & 0 & 1 & 0\end{pmatrix}.
\]
\end{definition}

Sa matrice échange les composantes $10$ et $11$ :
$\mathrm{CNOT}\ket{00}=\ket{00}$, $\mathrm{CNOT}\ket{01}=\ket{01}$,
$\mathrm{CNOT}\ket{10}=\ket{11}$, $\mathrm{CNOT}\ket{11}=\ket{10}$, soit
$\ket{i,j}\to\ket{i,i\oplus j}$ (addition modulo $2$).

On applique $H\otimes I$ puis $\mathrm{CNOT}$ à $\ket{00}$ (figure~\ref{fig:bell}) :
\[
  \ket{00}\ \xrightarrow{\ H\otimes I\ }\ \tfrac{1}{\sqrt2}\begin{pmatrix}1\\ 0\\ 1\\ 0\end{pmatrix}
  =\tfrac{1}{\sqrt2}\big(\ket{00}+\ket{10}\big)
  \ \xrightarrow{\ \mathrm{CNOT}\ }\ \tfrac{1}{\sqrt2}\begin{pmatrix}1\\ 0\\ 0\\ 1\end{pmatrix}
  =\tfrac{1}{\sqrt2}\big(\ket{00}+\ket{11}\big)=\ket{\Phi^+}.
\]

\begin{figure}[!ht]
\centering
\begin{tikzpicture}[>=Stealth, font=\small]
  % fils : B en haut (dernier chiffre du ket), A en bas (premier chiffre)
  \draw (0,1) -- (8.6,1);
  \draw (0,0) -- (8.6,0);
  \node[anchor=east] at (0,1) {$\ket{0}_B$};
  \node[anchor=east] at (0,0) {$\ket{0}_A$};
  % Hadamard sur A
  \node[draw=insarouge, fill=insarouge!8, minimum size=7mm, thick] at (2.2,0) {$H$};
  % CNOT : contrôle sur A (bas), cible sur B (haut)
  \draw[insarouge, thick] (5.4,0) -- (5.4,1.3);
  \fill[insarouge] (5.4,0) circle (2.6pt);
  \draw[insarouge, thick, fill=white] (5.4,1) circle (0.3);
  \draw[insarouge, thick] (5.1,1) -- (5.7,1);
  \draw[insarouge, thick] (5.4,0.7) -- (5.4,1.3);
  % accolade de sortie
  \draw[insagris, thick] (8.75,1.2) -- (8.9,1.2) -- (8.9,-0.2) -- (8.75,-0.2);
  \node[anchor=west] at (8.9,0.5) {$\ket{\Phi^+}_{AB}$};
  % états intermédiaires
  \foreach \x in {0.6,3.8,7.2}
    \draw[dashed, insagris!70] (\x,1.35) -- (\x,-0.95);
  \node[anchor=north, font=\footnotesize] at (0.6,-0.95) {$\ket{00}$};
  \node[anchor=north, font=\footnotesize] at (3.8,-0.95) {$\tfrac{1}{\sqrt2}(\ket{00}+\ket{10})$};
  \node[anchor=north, font=\footnotesize] at (7.2,-0.95) {$\tfrac{1}{\sqrt2}(\ket{00}+\ket{11})$};
\end{tikzpicture}
\caption{Circuit qui produit l'état de Bell $\ket{\Phi^+}$ : une porte de Hadamard sur $A$,
puis un CNOT ($A$ contrôle, $B$ cible). Les états intermédiaires sont indiqués sous les
traits pointillés. Convention : les fils sont dessinés dans l'ordre inverse de l'écriture
du ket $\ket{ij}_{AB}$, si bien que le premier chiffre (celui de $A$) est sur le fil du bas.}
\label{fig:bell}
\end{figure}

Les quatre états de base d'entrée donnent les quatre états de Bell :
\begin{center}
\renewcommand{\arraystretch}{1.6}
\begin{tabular}{@{}cccc@{}}
\toprule
Entrée & après $H\otimes I$ & après CNOT & Sortie \\
\midrule
$\ket{00}$ & $\frac{1}{\sqrt2}(\ket{00}+\ket{10})$ & $\frac{1}{\sqrt2}(\ket{00}+\ket{11})$ & $\ket{\Phi^+}$ \\
$\ket{01}$ & $\frac{1}{\sqrt2}(\ket{01}+\ket{11})$ & $\frac{1}{\sqrt2}(\ket{01}+\ket{10})$ & $\ket{\Psi^+}$ \\
$\ket{10}$ & $\frac{1}{\sqrt2}(\ket{00}-\ket{10})$ & $\frac{1}{\sqrt2}(\ket{00}-\ket{11})$ & $\ket{\Phi^-}$ \\
$\ket{11}$ & $\frac{1}{\sqrt2}(\ket{01}-\ket{11})$ & $\frac{1}{\sqrt2}(\ket{01}-\ket{10})$ & $\ket{\Psi^-}$ \\
\bottomrule
\end{tabular}
\end{center}
Sous forme matricielle, le produit $\mathrm{CNOT}\cdot(H\otimes I)$ échange les lignes $3$ et $4$ de
$H\otimes I$ :
\[
  \mathrm{CNOT}\,(H\otimes I)=\frac{1}{\sqrt2}\begin{pmatrix}1 & 0 & 1 & 0\\ 0 & 1 & 0 & 1\\ 0 & 1 & 0 & -1\\ 1 & 0 & -1 & 0\end{pmatrix},
\]
et ses quatre colonnes sont $\ket{\Phi^+}$, $\ket{\Psi^+}$, $\ket{\Phi^-}$, $\ket{\Psi^-}$.

\begin{remarque}[Porte intriquante]
Les états d'entrée $\ket{ij}$ sont séparables ($D=0$) ; les sorties sont intriquées
($D=\pm\frac12$). Aucune porte agissant séparément sur $A$ et sur $B$ ($U_A\otimes U_B$)
ne peut créer de l'intrication à partir d'un produit, puisque
$(U_A\otimes U_B)(\ket{u}\otimes\ket{v})=U_A\ket{u}\otimes U_B\ket{v}$ reste un produit.
Une porte à deux qubits comme le CNOT est indispensable.
\end{remarque}

% ---------------------------------------------------------------------
\subsection{Récapitulatif du chapitre}
% ---------------------------------------------------------------------

\begin{center}
\renewcommand{\arraystretch}{1.4}
\begin{tabular}{@{}ll@{}}
\toprule
\textbf{Notion} & \textbf{Formule} \\
\midrule
Espace de deux qubits & $\mathcal{H}^2\otimes\mathcal{H}^2=\mathcal{H}^4$, \quad dimension $2^n$ pour $n$ qubits \\
Produit de vecteurs & $(a_1,a_2)^{\mathsf T}\otimes(b_1,b_2)^{\mathsf T}=(a_1b_1,a_1b_2,a_2b_1,a_2b_2)^{\mathsf T}$ \\
Base de calcul & $\ket{ij}_{AB}=\ket{i}_A\otimes\ket{j}_B$, \quad $\braket{ij}{kl}=\delta_{ik}\delta_{jl}$ \\
État général & $\ket{\Psi}=\sum c_{ij}\ket{ij}$, \quad $\sum|c_{ij}|^2=1$, \quad $P(ij)=|c_{ij}|^2$ \\
Séparable & $\ket{\Psi}_A\otimes\ket{\Psi}_B$ \quad $\Leftrightarrow$ \quad $c_{00}c_{11}-c_{01}c_{10}=0$ \\
Intriqué & non séparable, \quad par exemple $\ket{\Phi^+}=\frac{1}{\sqrt2}(\ket{00}+\ket{11})$ \\
Matrices & $A\otimes B=(A_{ij}B)$, \quad $(A\otimes B)(\ket{u}\otimes\ket{v})=A\ket{u}\otimes B\ket{v}$ \\
CNOT & $\ket{i,j}\to\ket{i,i\oplus j}$, \quad $\mathrm{CNOT}\,(H\otimes I)\ket{00}=\ket{\Phi^+}$ \\
\bottomrule
\end{tabular}
\end{center}
.
%  Comme chapitre4.tex et chapitre5.tex, il ouvre une NOUVELLE SECTION
%  (numéro 4). Pour le rattacher à la section précédente, remplacer
%  \section par \subsection et chaque \subsection par \subsubsection.
%  Pas de préambule : environnements (definition, resultat, remarque),
%  macros (\ii, \ee, \dd, \ket, \bra, \braket) et couleurs (insarouge,
%  insagris) sont définis dans main.tex.
% =====================================================================

% Macros de Dirac (déjà définies dans main.tex récent ; ce bloc évite l'erreur
% « Undefined control sequence \ket » si main.tex est une ancienne version).
\providecommand{\ket}[1]{|#1\rangle}
\providecommand{\bra}[1]{\langle #1|}
\providecommand{\braket}[2]{\langle #1|#2\rangle}

% =====================================================================
\section{Systèmes multi-qubits et intrication}\label{sec:multi}
% =====================================================================

Un ordinateur quantique utilise plusieurs qubits. Cette section décrit l'espace des états
de deux qubits (le produit tensoriel), la mesure, la distinction entre états séparables
et états intriqués, puis les opérateurs qui agissent sur deux qubits. Tous les calculs
sont détaillés ; les qubits sont notés $A$ et $B$.

% ---------------------------------------------------------------------
\subsection{L'espace de deux qubits}\label{sec:tenseur}
% ---------------------------------------------------------------------

Chaque qubit vit dans $\mathcal{H}^2=\mathbb{C}^2$. L'espace du système $(A,B)$ est le
\emph{produit tensoriel}
\[
  \mathcal{H}^2\otimes\mathcal{H}^2=\mathcal{H}^4=\mathbb{C}^4 ,
\]
de dimension $2\times2=4$. Pour $n$ qubits la dimension est $2^n$ :

\begin{center}
\renewcommand{\arraystretch}{1.3}
\begin{tabular}{@{}lccccc@{}}
\toprule
Nombre de qubits $n$ & $1$ & $2$ & $3$ & $10$ & $50$ \\
\midrule
Dimension $2^n$ & $2$ & $4$ & $8$ & $1024$ & $\approx1{,}1\times10^{15}$ \\
\bottomrule
\end{tabular}
\end{center}

\begin{definition}{Produit tensoriel de deux vecteurs}
\[
  \begin{pmatrix}a_1\\ a_2\end{pmatrix}\otimes\begin{pmatrix}b_1\\ b_2\end{pmatrix}
  =\begin{pmatrix}a_1\begin{pmatrix}b_1\\ b_2\end{pmatrix}\\[2ex]
                  a_2\begin{pmatrix}b_1\\ b_2\end{pmatrix}\end{pmatrix}
  =\begin{pmatrix}a_1b_1\\ a_1b_2\\ a_2b_1\\ a_2b_2\end{pmatrix}.
\]
Chaque composante du premier vecteur multiplie le second vecteur en entier. On note
$\ket{ij}_{AB}=\ket{i}_A\otimes\ket{j}_B$ : le premier chiffre est l'état de $A$, le second
celui de $B$.
\end{definition}

\subsubsection{Les quatre états de base}

On développe chacun des quatre produits avec $\ket{0}=\begin{pmatrix}1\\ 0\end{pmatrix}$
et $\ket{1}=\begin{pmatrix}0\\ 1\end{pmatrix}$ :
\begin{align*}
  \ket{0}\otimes\ket{0}&=\begin{pmatrix}1\\ 0\end{pmatrix}\otimes\begin{pmatrix}1\\ 0\end{pmatrix}
    =\begin{pmatrix}1\cdot\begin{pmatrix}1\\ 0\end{pmatrix}\\[1.5ex]0\cdot\begin{pmatrix}1\\ 0\end{pmatrix}\end{pmatrix}
    =\begin{pmatrix}1\\ 0\\ 0\\ 0\end{pmatrix}=\ket{00},\\[1.5ex]
  \ket{0}\otimes\ket{1}&=\begin{pmatrix}1\\ 0\end{pmatrix}\otimes\begin{pmatrix}0\\ 1\end{pmatrix}
    =\begin{pmatrix}1\cdot\begin{pmatrix}0\\ 1\end{pmatrix}\\[1.5ex]0\cdot\begin{pmatrix}0\\ 1\end{pmatrix}\end{pmatrix}
    =\begin{pmatrix}0\\ 1\\ 0\\ 0\end{pmatrix}=\ket{01},\\[1.5ex]
  \ket{1}\otimes\ket{0}&=\begin{pmatrix}0\\ 1\end{pmatrix}\otimes\begin{pmatrix}1\\ 0\end{pmatrix}
    =\begin{pmatrix}0\cdot\begin{pmatrix}1\\ 0\end{pmatrix}\\[1.5ex]1\cdot\begin{pmatrix}1\\ 0\end{pmatrix}\end{pmatrix}
    =\begin{pmatrix}0\\ 0\\ 1\\ 0\end{pmatrix}=\ket{10},\\[1.5ex]
  \ket{1}\otimes\ket{1}&=\begin{pmatrix}0\\ 1\end{pmatrix}\otimes\begin{pmatrix}0\\ 1\end{pmatrix}
    =\begin{pmatrix}0\cdot\begin{pmatrix}0\\ 1\end{pmatrix}\\[1.5ex]1\cdot\begin{pmatrix}0\\ 1\end{pmatrix}\end{pmatrix}
    =\begin{pmatrix}0\\ 0\\ 0\\ 1\end{pmatrix}=\ket{11}.
\end{align*}
Le vecteur $\ket{ij}$ a son $1$ en position $2i+j$ (en comptant à partir de $0$) : c'est
la lecture de $ij$ en binaire. Les quatre vecteurs forment la base de calcul de
$\mathcal{H}^4$.

\subsubsection{Produit scalaire et orthonormalité}

\begin{resultat}{Produit scalaire d'un produit tensoriel}
\[
  \big(\bra{a}\otimes\bra{b}\big)\big(\ket{c}\otimes\ket{d}\big)=\braket{a}{c}\,\braket{b}{d}.
\]
\end{resultat}

\emph{Démonstration.} Les composantes de $\ket{a}\otimes\ket{b}$ sont $a_ib_j$, celles de
$\ket{c}\otimes\ket{d}$ sont $c_id_j$. Le produit scalaire est
$\sum_{i,j}(a_ib_j)^*c_id_j=\big(\sum_ia_i^*c_i\big)\big(\sum_jb_j^*d_j\big)=\braket{a}{c}\braket{b}{d}$.
\hfill$\square$

On en déduit $\braket{ij}{kl}=\braket{i}{k}\braket{j}{l}$. Les seize produits scalaires
des états de base sont :
\begin{center}
\renewcommand{\arraystretch}{1.3}
\begin{tabular}{@{}c|cccc@{}}
 & $\ket{00}$ & $\ket{01}$ & $\ket{10}$ & $\ket{11}$ \\
\hline
$\bra{00}$ & $1\cdot1=1$ & $1\cdot0=0$ & $0\cdot1=0$ & $0\cdot0=0$ \\
$\bra{01}$ & $1\cdot0=0$ & $1\cdot1=1$ & $0\cdot0=0$ & $0\cdot1=0$ \\
$\bra{10}$ & $0\cdot1=0$ & $0\cdot0=0$ & $1\cdot1=1$ & $1\cdot0=0$ \\
$\bra{11}$ & $0\cdot0=0$ & $0\cdot1=0$ & $1\cdot0=0$ & $1\cdot1=1$ \\
\end{tabular}
\end{center}
(chaque case est $\braket{i}{k}\cdot\braket{j}{l}$). On lit $\braket{ij}{kl}=\delta_{ik}\delta_{jl}$ :
la base est \emph{orthonormée}.

\subsubsection{Produit tensoriel d'états quelconques}

Le produit tensoriel est linéaire en chaque facteur :
$(\lambda\ket{a}+\mu\ket{a'})\otimes\ket{b}=\lambda\ket{a}\otimes\ket{b}+\mu\ket{a'}\otimes\ket{b}$,
et de même pour le second facteur. Pour $\alpha\ket{0}+\beta\ket{1}$ et
$\gamma\ket{0}+\delta\ket{1}$ :
\begin{align*}
  \big(\alpha\ket{0}+\beta\ket{1}\big)\otimes\big(\gamma\ket{0}+\delta\ket{1}\big)
  &=\alpha\ket{0}\otimes\big(\gamma\ket{0}+\delta\ket{1}\big)
    +\beta\ket{1}\otimes\big(\gamma\ket{0}+\delta\ket{1}\big)\\
  &=\alpha\gamma\ket{00}+\alpha\delta\ket{01}+\beta\gamma\ket{10}+\beta\delta\ket{11}.
\end{align*}
Par les composantes, $\begin{pmatrix}\alpha\\ \beta\end{pmatrix}\otimes\begin{pmatrix}\gamma\\ \delta\end{pmatrix}
=\begin{pmatrix}\alpha\gamma\\ \alpha\delta\\ \beta\gamma\\ \beta\delta\end{pmatrix}$ : c'est le même
résultat. La norme est conservée :
\[
  |\alpha\gamma|^2+|\alpha\delta|^2+|\beta\gamma|^2+|\beta\delta|^2
  =\big(|\alpha|^2+|\beta|^2\big)\big(|\gamma|^2+|\delta|^2\big)=1\cdot1=1 .
\]

\paragraph{Exemples.}
\begin{align*}
  \ket{+}\otimes\ket{+}&=\tfrac12\big(\ket{0}+\ket{1}\big)\otimes\big(\ket{0}+\ket{1}\big)
    =\tfrac12\big(\ket{00}+\ket{01}+\ket{10}+\ket{11}\big)
    =\tfrac12\begin{pmatrix}1\\ 1\\ 1\\ 1\end{pmatrix},\\[1ex]
  \ket{+}\otimes\ket{-}&=\tfrac12\big(\ket{0}+\ket{1}\big)\otimes\big(\ket{0}-\ket{1}\big)
    =\tfrac12\big(\ket{00}-\ket{01}+\ket{10}-\ket{11}\big)
    =\tfrac12\begin{pmatrix}1\\ -1\\ 1\\ -1\end{pmatrix},\\[1ex]
  \ket{0}\otimes\ket{+}&=\tfrac{1}{\sqrt2}\big(\ket{00}+\ket{01}\big)
    =\tfrac{1}{\sqrt2}\begin{pmatrix}1\\ 1\\ 0\\ 0\end{pmatrix},\qquad
  \ket{+}\otimes\ket{0}=\tfrac{1}{\sqrt2}\big(\ket{00}+\ket{10}\big)
    =\tfrac{1}{\sqrt2}\begin{pmatrix}1\\ 0\\ 1\\ 0\end{pmatrix}.
\end{align*}
Les deux derniers vecteurs sont différents : \emph{l'ordre des facteurs compte}, le premier
est toujours celui de $A$.

% ---------------------------------------------------------------------
\subsection{États à deux qubits et mesure}\label{sec:mesure2}
% ---------------------------------------------------------------------

Un état quelconque de $\mathcal{H}^4$ se décompose sur la base de calcul :
\[
  \ket{\Psi}=c_{00}\ket{00}+c_{01}\ket{01}+c_{10}\ket{10}+c_{11}\ket{11}
  =\begin{pmatrix}c_{00}\\ c_{01}\\ c_{10}\\ c_{11}\end{pmatrix},
  \qquad \sum_{i,j}|c_{ij}|^2=1 .
\]

\begin{resultat}{Mesure des deux qubits dans la base de calcul}
\[
  P(ij)=|\braket{ij}{\Psi}|^2=|c_{ij}|^2 .
\]
Pour un seul qubit : $P(A{=}0)=P(00)+P(01)$, $P(A{=}1)=P(10)+P(11)$,
$P(B{=}0)=P(00)+P(10)$, $P(B{=}1)=P(01)+P(11)$.
\end{resultat}

En effet $\braket{ij}{\Psi}=\sum_{k,l}c_{kl}\braket{ij}{kl}=\sum_{k,l}c_{kl}\delta_{ik}\delta_{jl}=c_{ij}$.

\paragraph{Exemple 1.} $\ket{\Psi}=\frac{1}{\sqrt2}\ket{00}+\frac{1}{\sqrt2}\ket{01}
=\big(\tfrac{1}{\sqrt2},\tfrac{1}{\sqrt2},0,0\big)^{\mathsf T}$. Il est normalisé :
$\frac12+\frac12+0+0=1$. Les probabilités jointes sont
$P(00)=P(01)=\frac12$ et $P(10)=P(11)=0$. Pour chaque qubit :
\[
  P(A{=}0)=\tfrac12+\tfrac12=1,\quad P(A{=}1)=0+0=0,\quad
  P(B{=}0)=\tfrac12+0=\tfrac12,\quad P(B{=}1)=\tfrac12+0=\tfrac12 .
\]
$A$ donne toujours $0$ ; $B$ donne $0$ ou $1$ au hasard.

\paragraph{Exemple 2 (mesure de $A$ seul).} Soit
$\ket{\Psi}=\frac12\ket{00}+\frac12\ket{01}+\frac{1}{\sqrt2}\ket{11}$. Normalisation :
$\frac14+\frac14+0+\frac12=1$. Probabilités jointes : $P(00)=\frac14$, $P(01)=\frac14$,
$P(10)=0$, $P(11)=\frac12$. Donc
\[
  P(A{=}0)=\tfrac14+\tfrac14=\tfrac12,\quad P(A{=}1)=0+\tfrac12=\tfrac12,\quad
  P(B{=}0)=\tfrac14+0=\tfrac14,\quad P(B{=}1)=\tfrac14+\tfrac12=\tfrac34 .
\]

\begin{remarque}[Effondrement après la mesure de $A$]
Si $A$ donne le résultat $a$, on ne garde que les termes de $\ket{\Psi}$ dont le premier
chiffre vaut $a$, puis on renormalise :
\[
  \ket{\Psi}\ \to\ \frac{c_{a0}\ket{a0}+c_{a1}\ket{a1}}{\sqrt{|c_{a0}|^2+|c_{a1}|^2}},
  \qquad\text{probabilité }|c_{a0}|^2+|c_{a1}|^2 .
\]
\end{remarque}

Pour l'exemple 2 :
\begin{itemize}[nosep]
  \item $A=0$ (probabilité $\frac12$) : on garde $\frac12\ket{00}+\frac12\ket{01}$, de norme
        au carré $\frac12$ ; renormalisé, $\frac{1}{\sqrt2}\big(\ket{00}+\ket{01}\big)
        =\ket{0}\otimes\ket{+}$. Ensuite $B$ vaut $0$ ou $1$ avec la probabilité $\frac12$.
  \item $A=1$ (probabilité $\frac12$) : on garde $\frac{1}{\sqrt2}\ket{11}$, de norme au carré
        $\frac12$ ; renormalisé, $\ket{11}=\ket{1}\otimes\ket{1}$. Ensuite $B$ vaut $1$
        avec certitude.
\end{itemize}
Contrôle : $P(B{=}0)=\frac12\cdot\frac12+\frac12\cdot0=\frac14$, comme calculé plus haut.
Ici la mesure de $A$ \emph{modifie} l'état de $B$ ($\ket{+}$ ou $\ket{1}$) : c'est la
signature d'un état intriqué (section suivante).

% ---------------------------------------------------------------------
\subsection{États séparables et états intriqués}\label{sec:separable}
% ---------------------------------------------------------------------

\begin{definition}{Séparable ou intriqué}
Un état $\ket{\Psi}_{AB}$ est \emph{séparable} s'il s'écrit comme un produit
$\ket{\Psi}_{AB}=\ket{\Psi}_A\otimes\ket{\Psi}_B$. Sinon il est \emph{intriqué} : on ne
peut pas attribuer d'état propre à chaque qubit.
\end{definition}

\paragraph{Exemple séparable (exemple 1 ci-dessus).} Par linéarité du produit tensoriel :
\begin{align*}
  \ket{\Psi}&=\tfrac{1}{\sqrt2}\ket{00}_{AB}+\tfrac{1}{\sqrt2}\ket{01}_{AB}
    =\tfrac{1}{\sqrt2}\Big(\ket{0}_A\otimes\ket{0}_B+\ket{0}_A\otimes\ket{1}_B\Big)\\
  &=\ket{0}_A\otimes\Big(\tfrac{1}{\sqrt2}\big(\ket{0}+\ket{1}\big)\Big)_B
    =\ket{0}_A\otimes\ket{+}_B ,
\end{align*}
donc $\ket{\Psi}_A=\ket{0}$ et $\ket{\Psi}_B=\ket{+}$. Contrôle par les composantes :
$\begin{pmatrix}1\\ 0\end{pmatrix}\otimes\tfrac{1}{\sqrt2}\begin{pmatrix}1\\ 1\end{pmatrix}
=\tfrac{1}{\sqrt2}\begin{pmatrix}1\cdot1\\ 1\cdot1\\ 0\cdot1\\ 0\cdot1\end{pmatrix}
=\tfrac{1}{\sqrt2}\begin{pmatrix}1\\ 1\\ 0\\ 0\end{pmatrix}$.

\subsubsection{Un critère de séparabilité}

Un produit quelconque a pour coefficients $c_{00}=ac$, $c_{01}=ad$, $c_{10}=bc$,
$c_{11}=bd$ (formule de la section~\ref{sec:tenseur}, avec $a,b$ pour $A$ et $c,d$ pour $B$), donc
\[
  c_{00}c_{11}=(ac)(bd)=abcd=(ad)(bc)=c_{01}c_{10}.
\]

\begin{resultat}{Test de séparabilité (deux qubits)}
$\ket{\Psi}=\sum c_{ij}\ket{ij}$ est séparable si et seulement si
\[
  D=c_{00}c_{11}-c_{01}c_{10}=0 .
\]
\end{resultat}

\emph{Démonstration.} ($\Rightarrow$) Calcul ci-dessus. ($\Leftarrow$) On regroupe les
termes selon l'état de $A$ :
\[
  \ket{\Psi}=\ket{0}\otimes\big(c_{00}\ket{0}+c_{01}\ket{1}\big)+\ket{1}\otimes\big(c_{10}\ket{0}+c_{11}\ket{1}\big).
\]
Si $D=0$, les lignes $(c_{00},c_{01})$ et $(c_{10},c_{11})$ sont proportionnelles. Si
$(c_{00},c_{01})\neq(0,0)$, il existe $\lambda$ avec $(c_{10},c_{11})=\lambda(c_{00},c_{01})$
et $\ket{\Psi}=\big(\ket{0}+\lambda\ket{1}\big)\otimes\big(c_{00}\ket{0}+c_{01}\ket{1}\big)$.
Si $(c_{00},c_{01})=(0,0)$, alors $\ket{\Psi}=\ket{1}\otimes\big(c_{10}\ket{0}+c_{11}\ket{1}\big)$.
Dans les deux cas $\ket{\Psi}$ est un produit ; on renormalise chaque facteur.
\hfill$\square$

\subsubsection{Application à plusieurs états}

\begin{center}
\renewcommand{\arraystretch}{1.5}
\begin{tabular}{@{}lccl@{}}
\toprule
\textbf{État} & $(c_{00},c_{01},c_{10},c_{11})$ & $D$ & \textbf{Nature} \\
\midrule
$\ket{00}$ & $(1,0,0,0)$ & $1\cdot0-0\cdot0=0$ & séparable \\
$\ket{+}\otimes\ket{+}$ & $\frac12(1,1,1,1)$ & $\frac14-\frac14=0$ & séparable \\
$\ket{+}\otimes\ket{-}$ & $\frac12(1,-1,1,-1)$ & $-\frac14-(-\frac14)=0$ & séparable \\
$\frac{1}{\sqrt2}(\ket{00}+\ket{01})$ & $\frac{1}{\sqrt2}(1,1,0,0)$ & $0-\frac12\cdot0=0$ & séparable \\
$\ket{\Phi^+}=\frac{1}{\sqrt2}(\ket{00}+\ket{11})$ & $\frac{1}{\sqrt2}(1,0,0,1)$ & $\frac12-0=\frac12$ & intriqué \\
$\ket{\Phi^-}=\frac{1}{\sqrt2}(\ket{00}-\ket{11})$ & $\frac{1}{\sqrt2}(1,0,0,-1)$ & $-\frac12-0=-\frac12$ & intriqué \\
$\ket{\Psi^+}=\frac{1}{\sqrt2}(\ket{01}+\ket{10})$ & $\frac{1}{\sqrt2}(0,1,1,0)$ & $0-\frac12=-\frac12$ & intriqué \\
$\ket{\Psi^-}=\frac{1}{\sqrt2}(\ket{01}-\ket{10})$ & $\frac{1}{\sqrt2}(0,1,-1,0)$ & $0-(-\frac12)=\frac12$ & intriqué \\
$\frac12\ket{00}+\frac12\ket{01}+\frac{1}{\sqrt2}\ket{11}$ & $(\frac12,\frac12,0,\frac{1}{\sqrt2})$ & $\frac{1}{2\sqrt2}-0$ & intriqué \\
\bottomrule
\end{tabular}
\end{center}

Les quatre états de Bell $\ket{\Phi^\pm}$, $\ket{\Psi^\pm}$ sont \emph{orthonormés} :
chacun est de norme $\frac12+\frac12=1$ ; $\ket{\Phi^\pm}$ n'ont des composantes que sur
$\ket{00}$ et $\ket{11}$, $\ket{\Psi^\pm}$ que sur $\ket{01}$ et $\ket{10}$, donc
$\braket{\Phi^\pm}{\Psi^\pm}=0$, et
$\braket{\Phi^+}{\Phi^-}=\frac12(1-1)=0$, $\braket{\Psi^+}{\Psi^-}=\frac12(1-1)=0$.
Ils forment une base de $\mathcal{H}^4$ constituée uniquement d'états intriqués.

\paragraph{Preuve directe que $\ket{\Phi^+}$ est intriqué.} Supposons
$\ket{\Phi^+}=(a\ket{0}+b\ket{1})\otimes(c\ket{0}+d\ket{1})$. En identifiant les
composantes : $ac=\frac{1}{\sqrt2}$, $ad=0$, $bc=0$, $bd=\frac{1}{\sqrt2}$. La condition
$ad=0$ impose $a=0$ (impossible car $ac\neq0$) ou $d=0$ (impossible car $bd\neq0$).
Contradiction : aucun produit ne convient.

\paragraph{Corrélations.} Pour un produit $\ket{\Psi}_A\otimes\ket{\Psi}_B$ avec
$\ket{\Psi}_A=a_0\ket{0}+a_1\ket{1}$ et $\ket{\Psi}_B=b_0\ket{0}+b_1\ket{1}$, on a
$P(ij)=|a_ib_j|^2=|a_i|^2|b_j|^2=P(A{=}i)\,P(B{=}j)$ : les deux qubits sont
\emph{indépendants}. Vérification avec $\ket{+}\otimes\ket{+}$ : $P(00)=\frac14=\frac12\cdot\frac12$.
Pour $\ket{\Phi^+}$ : $P(00)=P(11)=\frac12$ et $P(01)=P(10)=0$, donc
$P(A{=}0)=P(B{=}0)=\frac12$, mais
\[
  P(00)=\tfrac12\neq\tfrac12\cdot\tfrac12=P(A{=}0)\,P(B{=}0) .
\]
Les deux qubits donnent toujours le même résultat, aléatoire : ils sont corrélés.

\paragraph{État d'un qubit pris seul.} On prend la matrice de densité
$\rho=\ket{\Psi}\bra{\Psi}$ ($4\times4$, lignes et colonnes ordonnées $00,01,10,11$) et on
trace sur $B$ : $(\rho_A)_{ij}=\rho_{(i0),(j0)}+\rho_{(i1),(j1)}$.

\emph{État de Bell $\ket{\Phi^+}$.}
\[
  \rho=\frac12\begin{pmatrix}1 & 0 & 0 & 1\\ 0 & 0 & 0 & 0\\ 0 & 0 & 0 & 0\\ 1 & 0 & 0 & 1\end{pmatrix},
  \qquad
  \begin{aligned}
    (\rho_A)_{00}&=\rho_{00,00}+\rho_{01,01}=\tfrac12+0=\tfrac12,\\
    (\rho_A)_{01}&=\rho_{00,10}+\rho_{01,11}=0+0=0,\\
    (\rho_A)_{11}&=\rho_{10,10}+\rho_{11,11}=0+\tfrac12=\tfrac12 .
  \end{aligned}
\]
Donc $\rho_A=\frac12\mathrm{Id}$ : état \emph{mixte} ($\operatorname{Tr}\rho_A^2=\frac12$), au centre
de la boule de Bloch (section~\ref{sec:bloch}).

\emph{État séparable $\ket{0}\otimes\ket{+}$.}
\[
  \rho=\frac12\begin{pmatrix}1 & 1 & 0 & 0\\ 1 & 1 & 0 & 0\\ 0 & 0 & 0 & 0\\ 0 & 0 & 0 & 0\end{pmatrix},
  \qquad
  \begin{aligned}
    (\rho_A)_{00}&=\rho_{00,00}+\rho_{01,01}=\tfrac12+\tfrac12=1,\\
    (\rho_A)_{01}&=\rho_{00,10}+\rho_{01,11}=0+0=0,\\
    (\rho_A)_{11}&=\rho_{10,10}+\rho_{11,11}=0+0=0 .
  \end{aligned}
\]
Donc $\rho_A=\ket{0}\bra{0}$ : état \emph{pur}, celui du facteur $\ket{\Psi}_A$.

Un qubit d'un état séparable garde son propre état pur ; un qubit d'un état de Bell est
complètement aléatoire alors que le couple est parfaitement défini.

% ---------------------------------------------------------------------
\subsection{Opérateurs sur deux qubits}\label{sec:op2}
% ---------------------------------------------------------------------

\begin{definition}{Produit tensoriel de deux matrices}
\[
  A\otimes B=\begin{pmatrix}A_{11} & A_{12}\\ A_{21} & A_{22}\end{pmatrix}\otimes B
  =\begin{pmatrix}A_{11}B & A_{12}B\\ A_{21}B & A_{22}B\end{pmatrix}
  \qquad(\text{matrice }4\times4).
\]
\end{definition}

\begin{resultat}{Action sur un produit}
\[
  (A\otimes B)\big(\ket{u}\otimes\ket{v}\big)=A\ket{u}\otimes B\ket{v}.
\]
\end{resultat}

\emph{Démonstration.} Avec $\ket{u}=(u_1,u_2)^{\mathsf T}$, le vecteur $\ket{u}\otimes\ket{v}$ a
pour blocs $u_k\ket{v}$ ($k=1,2$). Le bloc de lignes $i$ de $(A\otimes B)(\ket{u}\otimes\ket{v})$ est
donc
\[
  \sum_kA_{ik}B\,\big(u_k\ket{v}\big)=\Big(\sum_kA_{ik}u_k\Big)B\ket{v}=(A\ket{u})_i\,B\ket{v},
\]
qui est le bloc $i$ de $A\ket{u}\otimes B\ket{v}$. \hfill$\square$

\subsubsection{Portes agissant sur un seul des deux qubits}

On note $I=\mathrm{Id}$ ($2\times2$). $X\otimes I$ agit sur $A$, $I\otimes X$ agit sur $B$ :
\begin{align*}
  I\otimes X&=\begin{pmatrix}1\cdot X & 0\cdot X\\ 0\cdot X & 1\cdot X\end{pmatrix}
    =\begin{pmatrix}0 & 1 & 0 & 0\\ 1 & 0 & 0 & 0\\ 0 & 0 & 0 & 1\\ 0 & 0 & 1 & 0\end{pmatrix},\\[1ex]
  X\otimes I&=\begin{pmatrix}0\cdot I & 1\cdot I\\ 1\cdot I & 0\cdot I\end{pmatrix}
    =\begin{pmatrix}0 & 0 & 1 & 0\\ 0 & 0 & 0 & 1\\ 1 & 0 & 0 & 0\\ 0 & 1 & 0 & 0\end{pmatrix}.
\end{align*}
Sur $\ket{00}=(1,0,0,0)^{\mathsf T}$ (première colonne de la matrice) :
$(I\otimes X)\ket{00}=(0,1,0,0)^{\mathsf T}=\ket{01}$ et
$(X\otimes I)\ket{00}=(0,0,1,0)^{\mathsf T}=\ket{10}$. Avec la règle :
$(I\otimes X)\ket{00}=I\ket{0}\otimes X\ket{0}=\ket{0}\otimes\ket{1}=\ket{01}$.

Pour la porte de Hadamard, avec $H=\frac{1}{\sqrt2}\begin{pmatrix}1 & 1\\ 1 & -1\end{pmatrix}$ :
\[
  H\otimes I=\frac{1}{\sqrt2}\begin{pmatrix}1\cdot I & 1\cdot I\\ 1\cdot I & -1\cdot I\end{pmatrix}
  =\frac{1}{\sqrt2}\begin{pmatrix}1 & 0 & 1 & 0\\ 0 & 1 & 0 & 1\\ 1 & 0 & -1 & 0\\ 0 & 1 & 0 & -1\end{pmatrix}.
\]
Ses colonnes donnent l'image des quatre états de base :
\begin{align*}
  (H\otimes I)\ket{00}&=\tfrac{1}{\sqrt2}(1,0,1,0)^{\mathsf T}=\ket{+}\otimes\ket{0},&
  (H\otimes I)\ket{01}&=\tfrac{1}{\sqrt2}(0,1,0,1)^{\mathsf T}=\ket{+}\otimes\ket{1},\\
  (H\otimes I)\ket{10}&=\tfrac{1}{\sqrt2}(1,0,-1,0)^{\mathsf T}=\ket{-}\otimes\ket{0},&
  (H\otimes I)\ket{11}&=\tfrac{1}{\sqrt2}(0,1,0,-1)^{\mathsf T}=\ket{-}\otimes\ket{1}.
\end{align*}

\subsubsection{Hadamard sur les deux qubits}

\[
  H\otimes H=\frac12\begin{pmatrix}1\cdot H' & 1\cdot H'\\ 1\cdot H' & -1\cdot H'\end{pmatrix},
  \quad H'=\begin{pmatrix}1 & 1\\ 1 & -1\end{pmatrix},
  \qquad
  H\otimes H=\frac12\begin{pmatrix}1 & 1 & 1 & 1\\ 1 & -1 & 1 & -1\\ 1 & 1 & -1 & -1\\ 1 & -1 & -1 & 1\end{pmatrix}.
\]
La première colonne donne $(H\otimes H)\ket{00}=\frac12(1,1,1,1)^{\mathsf T}=\ket{+}\otimes\ket{+}$
(cohérent avec $H\ket{0}\otimes H\ket{0}$) : les quatre résultats $00,01,10,11$ sont
équiprobables, chacun avec la probabilité $\frac14$. Sur $n$ qubits, $H^{\otimes n}\ket{0\ldots0}$
donne une superposition uniforme de $2^n$ états.

\subsubsection{La porte CNOT et les états de Bell}

\begin{definition}{Porte CNOT (NOT contrôlé)}
Le qubit $A$ est le \emph{contrôle}, $B$ la \emph{cible} : $B$ est inversé si $A=1$.
\[
  \mathrm{CNOT}=\ket{0}\bra{0}\otimes I+\ket{1}\bra{1}\otimes X
  =\begin{pmatrix}I & 0\\ 0 & X\end{pmatrix}
  =\begin{pmatrix}1 & 0 & 0 & 0\\ 0 & 1 & 0 & 0\\ 0 & 0 & 0 & 1\\ 0 & 0 & 1 & 0\end{pmatrix}.
\]
\end{definition}

Sa matrice échange les composantes $10$ et $11$ :
$\mathrm{CNOT}\ket{00}=\ket{00}$, $\mathrm{CNOT}\ket{01}=\ket{01}$,
$\mathrm{CNOT}\ket{10}=\ket{11}$, $\mathrm{CNOT}\ket{11}=\ket{10}$, soit
$\ket{i,j}\to\ket{i,i\oplus j}$ (addition modulo $2$).

On applique $H\otimes I$ puis $\mathrm{CNOT}$ à $\ket{00}$ (figure~\ref{fig:bell}) :
\[
  \ket{00}\ \xrightarrow{\ H\otimes I\ }\ \tfrac{1}{\sqrt2}\begin{pmatrix}1\\ 0\\ 1\\ 0\end{pmatrix}
  =\tfrac{1}{\sqrt2}\big(\ket{00}+\ket{10}\big)
  \ \xrightarrow{\ \mathrm{CNOT}\ }\ \tfrac{1}{\sqrt2}\begin{pmatrix}1\\ 0\\ 0\\ 1\end{pmatrix}
  =\tfrac{1}{\sqrt2}\big(\ket{00}+\ket{11}\big)=\ket{\Phi^+}.
\]

\begin{figure}[!ht]
\centering
\begin{tikzpicture}[>=Stealth, font=\small]
  % fils : B en haut (dernier chiffre du ket), A en bas (premier chiffre)
  \draw (0,1) -- (8.6,1);
  \draw (0,0) -- (8.6,0);
  \node[anchor=east] at (0,1) {$\ket{0}_B$};
  \node[anchor=east] at (0,0) {$\ket{0}_A$};
  % Hadamard sur A
  \node[draw=insarouge, fill=insarouge!8, minimum size=7mm, thick] at (2.2,0) {$H$};
  % CNOT : contrôle sur A (bas), cible sur B (haut)
  \draw[insarouge, thick] (5.4,0) -- (5.4,1.3);
  \fill[insarouge] (5.4,0) circle (2.6pt);
  \draw[insarouge, thick, fill=white] (5.4,1) circle (0.3);
  \draw[insarouge, thick] (5.1,1) -- (5.7,1);
  \draw[insarouge, thick] (5.4,0.7) -- (5.4,1.3);
  % accolade de sortie
  \draw[insagris, thick] (8.75,1.2) -- (8.9,1.2) -- (8.9,-0.2) -- (8.75,-0.2);
  \node[anchor=west] at (8.9,0.5) {$\ket{\Phi^+}_{AB}$};
  % états intermédiaires
  \foreach \x in {0.6,3.8,7.2}
    \draw[dashed, insagris!70] (\x,1.35) -- (\x,-0.95);
  \node[anchor=north, font=\footnotesize] at (0.6,-0.95) {$\ket{00}$};
  \node[anchor=north, font=\footnotesize] at (3.8,-0.95) {$\tfrac{1}{\sqrt2}(\ket{00}+\ket{10})$};
  \node[anchor=north, font=\footnotesize] at (7.2,-0.95) {$\tfrac{1}{\sqrt2}(\ket{00}+\ket{11})$};
\end{tikzpicture}
\caption{Circuit qui produit l'état de Bell $\ket{\Phi^+}$ : une porte de Hadamard sur $A$,
puis un CNOT ($A$ contrôle, $B$ cible). Les états intermédiaires sont indiqués sous les
traits pointillés. Convention : les fils sont dessinés dans l'ordre inverse de l'écriture
du ket $\ket{ij}_{AB}$, si bien que le premier chiffre (celui de $A$) est sur le fil du bas.}
\label{fig:bell}
\end{figure}

Les quatre états de base d'entrée donnent les quatre états de Bell :
\begin{center}
\renewcommand{\arraystretch}{1.6}
\begin{tabular}{@{}cccc@{}}
\toprule
Entrée & après $H\otimes I$ & après CNOT & Sortie \\
\midrule
$\ket{00}$ & $\frac{1}{\sqrt2}(\ket{00}+\ket{10})$ & $\frac{1}{\sqrt2}(\ket{00}+\ket{11})$ & $\ket{\Phi^+}$ \\
$\ket{01}$ & $\frac{1}{\sqrt2}(\ket{01}+\ket{11})$ & $\frac{1}{\sqrt2}(\ket{01}+\ket{10})$ & $\ket{\Psi^+}$ \\
$\ket{10}$ & $\frac{1}{\sqrt2}(\ket{00}-\ket{10})$ & $\frac{1}{\sqrt2}(\ket{00}-\ket{11})$ & $\ket{\Phi^-}$ \\
$\ket{11}$ & $\frac{1}{\sqrt2}(\ket{01}-\ket{11})$ & $\frac{1}{\sqrt2}(\ket{01}-\ket{10})$ & $\ket{\Psi^-}$ \\
\bottomrule
\end{tabular}
\end{center}
Sous forme matricielle, le produit $\mathrm{CNOT}\cdot(H\otimes I)$ échange les lignes $3$ et $4$ de
$H\otimes I$ :
\[
  \mathrm{CNOT}\,(H\otimes I)=\frac{1}{\sqrt2}\begin{pmatrix}1 & 0 & 1 & 0\\ 0 & 1 & 0 & 1\\ 0 & 1 & 0 & -1\\ 1 & 0 & -1 & 0\end{pmatrix},
\]
et ses quatre colonnes sont $\ket{\Phi^+}$, $\ket{\Psi^+}$, $\ket{\Phi^-}$, $\ket{\Psi^-}$.

\begin{remarque}[Porte intriquante]
Les états d'entrée $\ket{ij}$ sont séparables ($D=0$) ; les sorties sont intriquées
($D=\pm\frac12$). Aucune porte agissant séparément sur $A$ et sur $B$ ($U_A\otimes U_B$)
ne peut créer de l'intrication à partir d'un produit, puisque
$(U_A\otimes U_B)(\ket{u}\otimes\ket{v})=U_A\ket{u}\otimes U_B\ket{v}$ reste un produit.
Une porte à deux qubits comme le CNOT est indispensable.
\end{remarque}

% ---------------------------------------------------------------------
\subsection{Récapitulatif du chapitre}
% ---------------------------------------------------------------------

\begin{center}
\renewcommand{\arraystretch}{1.4}
\begin{tabular}{@{}ll@{}}
\toprule
\textbf{Notion} & \textbf{Formule} \\
\midrule
Espace de deux qubits & $\mathcal{H}^2\otimes\mathcal{H}^2=\mathcal{H}^4$, \quad dimension $2^n$ pour $n$ qubits \\
Produit de vecteurs & $(a_1,a_2)^{\mathsf T}\otimes(b_1,b_2)^{\mathsf T}=(a_1b_1,a_1b_2,a_2b_1,a_2b_2)^{\mathsf T}$ \\
Base de calcul & $\ket{ij}_{AB}=\ket{i}_A\otimes\ket{j}_B$, \quad $\braket{ij}{kl}=\delta_{ik}\delta_{jl}$ \\
État général & $\ket{\Psi}=\sum c_{ij}\ket{ij}$, \quad $\sum|c_{ij}|^2=1$, \quad $P(ij)=|c_{ij}|^2$ \\
Séparable & $\ket{\Psi}_A\otimes\ket{\Psi}_B$ \quad $\Leftrightarrow$ \quad $c_{00}c_{11}-c_{01}c_{10}=0$ \\
Intriqué & non séparable, \quad par exemple $\ket{\Phi^+}=\frac{1}{\sqrt2}(\ket{00}+\ket{11})$ \\
Matrices & $A\otimes B=(A_{ij}B)$, \quad $(A\otimes B)(\ket{u}\otimes\ket{v})=A\ket{u}\otimes B\ket{v}$ \\
CNOT & $\ket{i,j}\to\ket{i,i\oplus j}$, \quad $\mathrm{CNOT}\,(H\otimes I)\ket{00}=\ket{\Phi^+}$ \\
\bottomrule
\end{tabular}
\end{center}
